% !TEX program = xelatex
%% =====================================================================
%%  BÁO CÁO BÀI TẬP LỚN - DỰ ÁN 3 (CO2037)
%%  Xây dựng tài liệu thực hành mới với LTspice
%%  Mẫu báo cáo theo "Phần B" của đề (TS. Lê Trọng Nhân), dựng bằng LaTeX
%%  theo bố cục của các file logbook.
%%
%%  Biên dịch (bắt buộc XeLaTeX, chạy 2 lần để mục lục đúng):
%%      xelatex Project3_Report.tex      hoặc      latexmk Project3_Report.tex
%%  Overleaf: Menu -> Compiler -> XeLaTeX
%%
%%  CÁCH DÙNG NHANH
%%   1. Điền thông tin nhóm ở khối "THÔNG TIN NHÓM" bên dưới.
%%   2. Làm theo thứ tự các mục; thay mỗi KHUNG VIỀN CAM bằng nội dung của nhóm.
%%   3. Hình: đặt ảnh vào thư mục images/ đúng tên gợi ý -> khung cam tự thành hình.
%%   4. Khi xong: đổi \guidetrue thành \guidefalse để ẩn toàn bộ hướng dẫn.
%% =====================================================================
\documentclass[12pt,a4paper]{article}

\usepackage[a4paper,top=2.4cm,bottom=2.5cm,left=2.5cm,right=2.5cm,headheight=15pt]{geometry}
\usepackage{fontspec}
\setmainfont{Latin Modern Roman}
\setsansfont{Latin Modern Sans}
\setmonofont{Latin Modern Mono}[AutoFakeBold=1.5]

\usepackage[table]{xcolor}
\usepackage{colortbl,array,tabularx,longtable,multirow}
\usepackage{enumitem,graphicx,float,caption,amsmath,amssymb,etoolbox}
\usepackage{tikz}
\usetikzlibrary{shapes.geometric,arrows.meta,positioning,calc}
\usepackage[most]{tcolorbox}
\usepackage{fancyhdr,titlesec}
\setcounter{secnumdepth}{2}\setcounter{tocdepth}{2}
\usepackage[colorlinks=true,linkcolor=navy,urlcolor=navy,citecolor=navy]{hyperref}

%% ===================== BẬT / TẮT HƯỚNG DẪN =====================
\newif\ifguide
\guidetrue          % <-- đổi thành \guidefalse khi nộp bài (ẩn mọi hướng dẫn)

%% ===================== THÔNG TIN NHÓM (sửa tại đây) =====================
\newcommand{\NhomSo}{9}                          % Nhóm
\newcommand{\LopHoc}{\dotuline}                  % Lớp (điền, ví dụ: CC04)
\newcommand{\HocKy}{\dotuline}                   % Học kỳ (điền)
\newcommand{\TVa}{Phạm Đỗ Hồng Phúc}   \newcommand{\MSa}{2510359}   \newcommand{\PTa}{-- Lab 5, Lab 6}  % lab phụ trách của TV1
\newcommand{\TVb}{Lê Lương Vỹ}          \newcommand{\MSb}{2214056}   \newcommand{\PTb}{-- Lab 1, Lab 2}  % lab phụ trách của TV2
\newcommand{\TVc}{Nguyễn Bảo Duy Thành} \newcommand{\MSc}{2510372}   \newcommand{\PTc}{-- Lab 3, Lab 4}  % lab phụ trách của TV3
\newcommand{\PhienBanLT}{LTspice 24.1.8 (64-bit) / Windows 11}                       % Phiên bản LTspice / hệ điều hành
\newcommand{\CongCuSoanThao}{LaTeX (XeLaTeX)}    % Công cụ biên soạn tài liệu
\newcommand{\CongCuAI}{Antigravity Agent}                         % Công cụ AI đã sử dụng
%% =========================================================================

\definecolor{navy}{HTML}{1F2F5E}
\definecolor{accent}{HTML}{D98200}
\definecolor{guidegray}{HTML}{6E6E6E}
\definecolor{lightnavy}{HTML}{EAEFF8}
\definecolor{lightorange}{HTML}{FFF3DC}
\newcommand{\dotuline}{\makebox[2.6cm]{\dotfill}}

%% ---------- Việt hoá tên bảng/hình, caption ----------
\renewcommand{\tablename}{Bảng}
\renewcommand{\figurename}{Hình}
\renewcommand{\contentsname}{MỤC LỤC}
\captionsetup{font={small,it},labelfont={small,bf,it},labelsep=colon,skip=5pt}

%% ---------- Trang, chân trang (giống logbook: đường kẻ + chữ nghiêng) ----------
\pagestyle{fancy}
\fancyhf{}
\fancyhead[R]{\small\itshape\color{guidegray} CO2037 -- Dự án 3: Tài liệu thực hành LTspice}
\fancyfoot[L]{\small\itshape\color{guidegray} HCMUT -- Khoa Khoa học và Kỹ thuật Máy tính}
\fancyfoot[R]{\small\itshape\color{guidegray} Trang \thepage}
\renewcommand{\headrulewidth}{0pt}
\renewcommand{\footrulewidth}{0.8pt}
\setlength{\parindent}{0pt}
\setlength{\parskip}{5pt}
\linespread{1.08}
\setlist{itemsep=2pt,topsep=3pt,parsep=0pt}

%% ---------- Tiêu đề mục (xanh navy, đậm) ----------
\titleformat{\section}{\sffamily\bfseries\Large\color{navy}}{\thesection.}{0.6em}{}
\titleformat{\subsection}{\sffamily\bfseries\large\color{navy}}{\thesubsection.}{0.6em}{}
\titleformat{\subsubsection}{\sffamily\bfseries\normalsize\color{navy}}{}{0pt}{}
\titlespacing*{\section}{0pt}{16pt}{6pt}
\titlespacing*{\subsection}{0pt}{12pt}{4pt}
\titlespacing*{\subsubsection}{0pt}{8pt}{3pt}

%% =====================================================================
%%  MACRO GIÚP LÀM BÀI
%% =====================================================================
% (1) Chữ xám nghiêng = hướng dẫn của đề -> xoá khi hoàn thành
\newcommand{\hd}[1]{\ifguide{\par\color{guidegray}\itshape #1\par}\fi}

% (2) Khung xanh = "Cách làm" chi tiết do mẫu này bổ sung
\newcommand{\cachlam}[2][Cách làm]{%
  \ifguide
  \begin{tcolorbox}[enhanced,breakable,colback=lightnavy,colframe=navy,boxrule=0.6pt,arc=2pt,
    left=6pt,right=6pt,top=3pt,bottom=3pt,fonttitle=\sffamily\bfseries\small,title={#1}]
  \small #2
  \end{tcolorbox}
  \fi}

% (3) Khung viền cam = vị trí chèn nội dung của nhóm
\newcommand{\chen}[2][2.2cm]{%
  \begin{tcolorbox}[colback=lightorange,colframe=accent,boxrule=1.1pt,arc=2pt,
     height=#1,valign=center,halign=center,before skip=6pt,after skip=6pt]
  \color{accent!85!black}\bfseries\itshape [#2]
  \end{tcolorbox}}

% (4) Chèn hình: nếu có file images/<tên> thì tự thành hình, nếu chưa có thì hiện khung cam
\newcommand{\hinhanh}[4]{% tên file, độ rộng, chú thích, nhãn
  \IfFileExists{images/#1}{%
    \begin{figure}[H]\centering\includegraphics[width=#2]{images/#1}%
      \caption{#3}\label{#4}\end{figure}}{%
    \chen[2.6cm]{Chèn hình: #3 \\[2pt] {\small\normalfont\upshape(đặt file \texttt{images/#1})}}}}

% (5) Ô tick và dấu
\newcommand{\cb}{$\square$\ }
\newcommand{\ar}{\texorpdfstring{$\rightarrow$}{->}}

% (6) Hàng tiêu đề bảng
\newcommand{\hc}[1]{\cellcolor{navy}\color{white}\sffamily\bfseries\small #1}
\newcommand{\hrulenavy}{\arrayrulecolor{navy}}
\arrayrulecolor{navy}

% (7) Bảng 10 bài tập của một lab
\newcommand{\labrow}[1]{#1 & & & & & \rule{0pt}{0.95cm}\\ \hline}
\newcommand{\labbang}[1]{%
  \begin{table}[H]\centering\small
  \caption{Danh sách bài tập Lab #1}\label{tab:lab#1}
  \setlength{\tabcolsep}{3.5pt}
  \begin{tabular}{|c|p{2.5cm}|p{3.3cm}|p{1.9cm}|p{2.7cm}|p{2.2cm}|}
  \hline
  \hc{\#} & \hc{Bài gốc} & \hc{Bài mới (LTspice)} & \hc{Trạng thái} & \hc{Lệnh mô phỏng} & \hc{Sai lệch LT -- MP}\\ \hline
  \labrow{1}\labrow{2}\labrow{3}\labrow{4}\labrow{5}\labrow{6}\labrow{7}\labrow{8}\labrow{9}\labrow{10}
  \end{tabular}
  \par\vspace{2pt}{\footnotesize Trạng thái: \textbf{Giữ nguyên} (chỉ đổi công cụ) / \textbf{Điều chỉnh} (đổi linh kiện, giá trị nhưng giữ mục tiêu) / \textbf{Thay thế} (bài tương đương mới).}
  \end{table}}

% (8) Bảng bài tập nâng cao của một lab
\newcommand{\advrow}[1]{#1 & & & \rule{0pt}{0.9cm}\\ \hline}
\newcommand{\advbang}[1]{%
  \begin{table}[H]\centering\small
  \caption{Bài tập nâng cao Lab #1 (tối thiểu 2 bài)}\label{tab:adv#1}
  \setlength{\tabcolsep}{4pt}
  \begin{tabular}{|p{1.5cm}|p{3.6cm}|p{4.2cm}|p{5.2cm}|}
  \hline
  \hc{Mã file} & \hc{Tên bài} & \hc{Mục tiêu} & \hc{Mô tả ngắn}\\ \hline
  \advrow{Adv01}\advrow{Adv02}\advrow{Adv03 (nếu có)}
  \end{tabular}
  \end{table}}

% (9) Một câu hỏi trong mục 7
\newcommand{\cau}[3]{% số câu, nội dung câu hỏi (trích từ đề), cách làm
  \subsubsection{Câu [#1]}
  \begin{tcolorbox}[colback=white,colframe=guidegray!70,boxrule=0.4pt,arc=1pt,left=5pt,right=5pt,top=2pt,bottom=2pt]
  \small\textit{Đề:} #2
  \end{tcolorbox}
  \cachlam[Cách làm câu {[#1]}]{#3}
  \chen[2.6cm]{Trả lời của nhóm}}

\newcommand{\cauhoi}[2]{% số câu, nội dung câu hỏi (trích từ đề)
  \subsubsection{Câu [#1]}
  \begin{tcolorbox}[colback=white,colframe=guidegray!70,boxrule=0.4pt,arc=1pt,left=5pt,right=5pt,top=2pt,bottom=2pt]
  \small\textit{Đề:} #2
  \end{tcolorbox}
  \par\vspace{4pt}}

%% =====================================================================
\begin{document}

%% ============================== BÌA ==============================
\begin{titlepage}
\centering
{\sffamily\bfseries\large ĐẠI HỌC QUỐC GIA TP. HỒ CHÍ MINH\\
TRƯỜNG ĐẠI HỌC BÁCH KHOA\\
KHOA KHOA HỌC VÀ KỸ THUẬT MÁY TÍNH\par}
\vspace{1.6cm}
{\sffamily\bfseries\Large\color{navy} LINH KIỆN VÀ MẠCH ĐIỆN TỬ (CO2037)\par}
{\itshape\color{navy} Electronic Devices and Circuits\par}
\vspace{1.8cm}
{\sffamily\bfseries\LARGE\color{accent} BÁO CÁO BÀI TẬP LỚN -- DỰ ÁN 3\par}
\vspace{4pt}
{\sffamily\bfseries\Huge\color{navy} XÂY DỰNG TÀI LIỆU\\[3pt] THỰC HÀNH MỚI\\[3pt] VỚI LTSPICE\par}
\vspace{6pt}
{\itshape\large Chuyển 6 bài thực hành mô phỏng từ PSpice for TI sang LTspice\par}
\vspace{1.6cm}
\renewcommand{\arraystretch}{1.5}
\begin{tabular}{|>{\bfseries\sffamily\small}p{4.4cm}|p{9.4cm}|}
\hline
\hc{Thông tin} & \hc{Nội dung}\\ \hline
Giảng viên & TS. Lê Trọng Nhân\\ \hline
Nhóm & \NhomSo \hspace{1.2cm} \textbf{Lớp:} \LopHoc\\ \hline
Học kỳ & \HocKy\\ \hline
Thành viên 1 & \TVa\ -- \MSa\\ \hline
Thành viên 2 & \TVb\ -- \MSb\\ \hline
Thành viên 3 & \TVc\ -- \MSc\\ \hline
\end{tabular}
\vfill
{\itshape TP. Hồ Chí Minh, ngày \dotuline\ tháng \dotuline\ năm \dotuline\par}
\end{titlepage}

\pagenumbering{arabic}
\tableofcontents
\clearpage

%% ================== PHẦN HƯỚNG DẪN THỰC HIỆN (ẩn khi \guidefalse) ==================
\ifguide
\section*{HƯỚNG DẪN THỰC HIỆN BÁO CÁO (đọc trước -- xoá khi nộp)}
\addcontentsline{toc}{section}{HƯỚNG DẪN THỰC HIỆN BÁO CÁO (đọc trước -- xoá khi nộp)}

\subsection*{0.1. Cách dùng file mẫu này}
\begin{itemize}
  \item \textbf{Chữ xám nghiêng} (như dòng này ở đầu mỗi mục): hướng dẫn của đề -- xoá khi hoàn thành.
  \item \textbf{Khung viền cam}: vị trí chèn nội dung của nhóm -- thay khung bằng nội dung thật.
  \item \textbf{Khung nền xanh} ``Cách làm'': hướng dẫn chi tiết cách thực hiện từng mục, bổ sung từ yêu cầu của đề.
  \item Ẩn toàn bộ hướng dẫn: đổi \texttt{\textbackslash guidetrue} thành \texttt{\textbackslash guidefalse} ở đầu file.
  \item Chèn hình: đặt ảnh vào thư mục \texttt{images/} đúng tên ghi trong khung cam; khi file có mặt, khung tự chuyển thành hình có đánh số và chú thích.
  \item Thông tin nhóm (họ tên, MSSV, lab phụ trách, phiên bản LTspice, công cụ AI): sửa ở khối \texttt{THÔNG TIN NHÓM} đầu file; các bảng bên dưới tự cập nhật.
\end{itemize}

\subsection*{0.2. Quy trình thực hiện dự án (Hình 1 của đề)}
\begin{center}
\resizebox{0.98\linewidth}{!}{%
\begin{tikzpicture}[
  node distance=0.5cm,
  box/.style={draw,rounded corners=4pt,very thick,align=center,minimum width=3.0cm,minimum height=1.6cm,font=\sffamily\small},
  arr/.style={-{Stealth[length=2.2mm]},thick,gray!70!black}]
  \node[box,draw=gray!70!black] (a) at (0,0) {Tài liệu cũ\\(PSpice for TI)\\6 lab $\times$ 10 bài};
  \node[box,draw=navy] (b) at (3.7,0) {Phân tích nội dung\\\& phân công\\(3 sinh viên)};
  \node[box,draw=accent] (c) at (7.4,0) {Cài đặt LTspice\\bản offline\\+ thư viện};
  \node[box,draw=navy] (d) at (11.1,0) {Hiện thực lại\\từng bài tập\\trên LTspice};
  \node[draw=accent,very thick,diamond,aspect=1.6,align=center,font=\sffamily\small,inner sep=1pt] (e) at (14.9,0) {Thực hiện\\được?};
  \node[box,draw=green!50!black,minimum width=2.7cm] (f) at (13.2,-3.0) {Giữ nguyên /\\điều chỉnh};
  \node[box,draw=accent,minimum width=2.7cm] (g) at (16.5,-3.0) {Đề xuất bài\\tương đương};
  \node[box,draw=green!50!black] (h) at (9.3,-5.4) {Thêm bài tập\\nâng cao\\(mỗi lab)};
  \node[box,draw=navy] (i) at (5.0,-5.4) {Kiểm thử chéo\\(thành viên khác\\làm theo tài liệu)};
  \node[box,draw=green!50!black] (j) at (0.8,-5.4) {Tài liệu mới\\+ lời giải\\+ file .asc};
  \draw[arr] (a)--(b); \draw[arr] (b)--(c); \draw[arr] (c)--(d); \draw[arr] (d)--(e);
  \draw[arr] (e)--node[left,green!50!black,font=\sffamily\small\bfseries]{Có}(f);
  \draw[arr] (e)--node[right,red!70!black,font=\sffamily\small\bfseries]{Không}(g);
  \draw[arr] (f)--(h.east); \draw[arr] (g)--(h.east);
  \draw[arr] (h)--(i); \draw[arr] (i)--(j);
\end{tikzpicture}}
\end{center}
\textbf{Thứ tự làm việc khuyến nghị:}
\begin{enumerate}
  \item \textbf{Phân tích và phân công:} đọc tài liệu gốc 6 chương; mỗi thành viên phụ trách chính \textbf{2 lab}; cả nhóm cùng làm cài đặt, bảng ánh xạ và kiểm thử chéo (mục 11).
  \item \textbf{Cài đặt offline:} tìm bộ cài, kiểm tra bằng một file mẫu trong \texttt{examples}, chụp màn hình (mục 2).
  \item \textbf{Hoàn thiện bảng ánh xạ thao tác} PSpice for TI $\rightarrow$ LTspice theo đúng phiên bản đang dùng (mục 3).
  \item \textbf{Hiện thực từng bài tập:} vẽ \texttt{.asc}, chạy mô phỏng, ghi lệnh mô phỏng, đối chiếu lý thuyết -- mô phỏng, gán trạng thái Giữ nguyên / Điều chỉnh / Thay thế (mục 5).
  \item \textbf{Bài không chuyển đổi được:} giải thích lý do và đề xuất bài tương đương (mục 6).
  \item \textbf{Thêm bài nâng cao:} tối thiểu 2 bài mỗi lab (mục 5).
  \item \textbf{Kiểm thử chéo:} người khác làm theo tài liệu, ghi lỗi, sửa tài liệu (mục 8).
  \item \textbf{Hoàn thiện} lời giải, báo cáo, nhật ký AI, phân công; đóng gói và nộp.
\end{enumerate}

\subsection*{0.3. Sản phẩm phải nộp}
\begin{table}[H]\centering\small
\caption{Danh sách sản phẩm nộp (theo đề)}
\renewcommand{\arraystretch}{1.25}
\begin{tabular}{|c|p{8.2cm}|p{4.2cm}|}
\hline
\hc{STT} & \hc{Sản phẩm} & \hc{Định dạng}\\ \hline
1 & Tài liệu thực hành mới (6 bài lab) & .pdf + file nguồn (.tex hoặc .docx)\\ \hline
2 & File lời giải tham khảo & .pdf\\ \hline
3 & Thư mục file mô phỏng LTspice và thư viện bổ sung & .zip\\ \hline
4 & Hướng dẫn cài đặt offline (có thể là một chương trong tài liệu) & .pdf\\ \hline
5 & Báo cáo dự án theo mẫu Phần B (chính là file này) & .docx và .pdf\\ \hline
6 & (Khuyến khích) Video demo ngắn cho mỗi lab & .mp4 hoặc liên kết\\ \hline
\end{tabular}
\end{table}
\textbf{Tên file nộp:} \texttt{MSSV1\_MSSV2\_MSSV3\_Project3.zip}. Tên file mô phỏng đặt thống nhất: \texttt{Lab1/Ex01.asc}, \dots, \texttt{Lab6/Adv02.asc}.

\cachlam[Lưu ý quan trọng về định dạng báo cáo]{
Đề yêu cầu báo cáo theo mẫu Phần B ở dạng \textbf{.docx và .pdf}. File LaTeX này cho ra bản \textbf{.pdf}; hãy \textbf{hỏi giảng viên} xem có chấp nhận file nguồn \texttt{.tex} thay cho \texttt{.docx} đối với báo cáo hay không. Nếu không, chuyển nội dung sang mẫu \texttt{.docx} của đề (cấu trúc 12 mục giống hệt file này).}

\subsection*{0.4. Những quy tắc bắt buộc của đề}
\begin{itemize}
  \item \cb Đủ \textbf{6 lab}, mỗi lab \textbf{đúng 10 bài tập chính} (chương gốc nhiều hơn thì chọn lọc/gộp, ít hơn thì bổ sung) và \textbf{$\geq$ 2 bài nâng cao}.
  \item \cb Mọi hình mạch và hình kết quả \textbf{chụp từ LTspice}, đủ rõ khi in, có đánh số và chú thích.
  \item \cb Mỗi bài có giá trị linh kiện cụ thể và ghi rõ mô hình linh kiện; mọi \texttt{.asc} \textbf{mở và chạy được ngay trên bộ cài offline} mà không cần tải thêm.
  \item \cb Phần tính toán \textbf{để trống} cho sinh viên; lời giải chỉ nằm trong file dành cho giảng viên.
  \item \cb Mọi số liệu do nhóm \textbf{chạy mô phỏng thật} trên LTspice -- \textbf{không dùng số liệu do AI tự sinh}.
  \item \cb Mọi lần dùng AI có ảnh hưởng tới kết quả phải ghi vào \textbf{Nhật ký sử dụng AI} (mục 10), kèm cách kiểm chứng.
  \item \cb Ngôn ngữ tài liệu thực hành: \textbf{tiếng Anh}, thuật ngữ và ký hiệu nhất quán giữa các bài.
  \item \cb Mọi thành viên phải giải thích được \textbf{mọi bài lab} (có thể vấn đáp), không chỉ phần mình phụ trách.
  \item \cb Ngoài phạm vi: Ch.~4 (Altium Designer), các mục Altium trong Ch.~5 và Ch.~6, Ch.~8 (Final Project); không thiết kế PCB, không thi công mạch.
\end{itemize}

\subsection*{0.5. Tiêu chí đánh giá và vị trí trong báo cáo này}
\begin{table}[H]\centering\small
\caption{Tiêu chí đánh giá Dự án 3 (theo đề) và nơi thể hiện}
\renewcommand{\arraystretch}{1.25}
\begin{tabular}{|p{4.2cm}|c|p{8.2cm}|}
\hline
\hc{Tiêu chí} & \hc{Tỉ trọng} & \hc{Thể hiện trong báo cáo này}\\ \hline
Tài liệu thực hành & 35\,\% & Nộp riêng; báo cáo mô tả ở mục 4, 5\\ \hline
File mô phỏng \& lời giải & 20\,\% & Nộp riêng; mô tả cấu trúc ở mục 9\\ \hline
Chuyển đổi \& bài tương đương & 15\,\% & Mục 3 (bảng ánh xạ), mục 5 (trạng thái), mục 6\\ \hline
Bài tập nâng cao & 10\,\% & Mục 5 (bảng nâng cao từng lab)\\ \hline
Cài đặt offline \& câu hỏi & 10\,\% & Mục 2 và mục 7 (câu [1]--[9])\\ \hline
Báo cáo \& quy trình & 10\,\% & Mục 8 (kiểm thử chéo), 10 (nhật ký AI), 11 (phân công)\\ \hline
\end{tabular}
\end{table}

\subsection*{0.6. Một bài tập chính phải có những gì}
Mỗi bài tập chính trong tài liệu thực hành mới gồm: (1) \textbf{sơ đồ mạch}; (2) \textbf{phần tính toán lý thuyết} (để trống cho sinh viên điền); (3) \textbf{phần mô phỏng} (hướng dẫn cấu hình lệnh); (4) \textbf{bảng so sánh} lý thuyết -- mô phỏng. Mỗi lab còn có phần \textbf{Introduction}, \textbf{hướng dẫn LTspice có ảnh chụp màn hình}, và bài nâng cao (ví dụ dùng \texttt{.step}, \texttt{.meas}, phân tích nhiệt độ, Monte Carlo, phân tích Fourier).

\clearpage
\fi

%% =====================================================================
%%                       NỘI DUNG BÁO CÁO (PHẦN B)
%% =====================================================================
\hd{Báo cáo này mô tả \textbf{quá trình} thực hiện dự án; bản thân tài liệu thực hành mới được nộp riêng. Các dòng chữ nghiêng màu xám là hướng dẫn -- xoá đi khi hoàn thành. Các khung viền cam là vị trí chèn nội dung.}

\begin{table}[H]\centering\small
\caption{Thông tin nhóm}
\renewcommand{\arraystretch}{1.5}
\begin{tabular}{|p{6.2cm}|p{8.8cm}|}
\hline
\hc{Thông tin} & \hc{Nội dung}\\ \hline
Thành viên 1 (Họ tên -- MSSV) -- lab phụ trách & \TVa\ -- \MSa\ \PTa\\ \hline
Thành viên 2 (Họ tên -- MSSV) -- lab phụ trách & \TVb\ -- \MSb\ \PTb\\ \hline
Thành viên 3 (Họ tên -- MSSV) -- lab phụ trách & \TVc\ -- \MSc\ \PTc\\ \hline
Lớp / Nhóm & \LopHoc\ / Nhóm \NhomSo\\ \hline
Phiên bản LTspice / hệ điều hành & \PhienBanLT\\ \hline
Công cụ biên soạn tài liệu (LaTeX/Word) & \CongCuSoanThao\\ \hline
Công cụ AI đã sử dụng & \CongCuAI\\ \hline
\end{tabular}
\end{table}
\cachlam{Điền \texttt{\textbackslash PTa}, \texttt{\textbackslash PTb}, \texttt{\textbackslash PTc} (lab phụ trách, ví dụ ``Lab 1, Lab 2''), \texttt{\textbackslash PhienBanLT} (ghi số phiên bản \textbf{đúng như hiển thị trong phần mềm} và hệ điều hành) và \texttt{\textbackslash CongCuAI} ở khối THÔNG TIN NHÓM đầu file. Lớp và học kỳ điền ở dòng \texttt{\textbackslash LopHoc}, \texttt{\textbackslash HocKy}.}

%% ----------------------------- 1 -----------------------------
\section{Giới thiệu}

\subsection{1.1. Bối cảnh và mục tiêu dự án}
Học phần \textit{Linh kiện và Mạch điện tử} (CO2037 / CO2104) là môn học kỹ thuật cơ sở cốt lõi trong chương trình đào tạo kỹ sư ngành Khoa học Máy tính và Kỹ thuật Máy tính tại Trường Đại học Bách Khoa -- ĐHQG-HCM. Học phần trang bị cho sinh viên kiến thức chuyên sâu về mô hình toán học, nguyên lý vật lý bán dẫn và kỹ năng phân tích thiết kế các mạch điện tử tương tự cơ bản (Diode, BJT, FET/MOSFET, Op-Amp, nguồn ổn áp và chuyển mạch công suất).

Trong nhiều năm qua, các bài thực hành mô phỏng được thực hiện trên phần mềm thương mại \textit{PSpice for TI} (Texas Instruments / Cadence). Mặc dù có ưu thế về hệ thống mô hình IC chính hãng của TI, phần mềm này bộc lộ nhiều rào cản kỹ thuật nghiêm trọng đối với sinh viên:
\begin{itemize}[itemsep=2pt]
  \item \textbf{Ràng buộc tài khoản và kích hoạt:} Bắt buộc phải đăng ký tài khoản TI, chờ phê duyệt mã cấp phép (access key) và phải kết nối Internet định kỳ để xác thực bản quyền, gây cản trở lớn khi sinh viên học tập tại phòng máy ngắt mạng hoặc khu vực không có kết nối ổn định.
  \item \textbf{Giới hạn tính năng khi dùng linh kiện ngoài TI:} PSpice for TI áp đặt giới hạn hiển thị tối đa 3 trace đồ thị và hạn chế bộ giải đối với các linh kiện không do TI sản xuất.
  \item \textbf{Cấu trúc quản lý dự án cồng kềnh:} Mỗi bài tập đòi hỏi tạo một project phức tạp với hàng loạt thư mục con và thiết lập Simulation Profile tách rời.
\end{itemize}

Trước thực trạng đó, \textbf{Mục tiêu trọng tâm của Dự án 3} là: \textbf{Nghiên cứu, chuyển đổi toàn diện và xây dựng bộ tài liệu thực hành mô phỏng mới dựa trên phần mềm Analog Devices LTspice} cho toàn bộ 6 bài Lab (Lab 1 đến Lab 6). Dự án hướng đến việc cung cấp một môi trường mô phỏng mở, miễn phí 100\%, hoạt động độc lập ngoại tuyến (offline), không giới hạn số nút mạch hay số linh kiện, giúp sinh viên làm chủ công cụ mô phỏng chuẩn công nghiệp ngay trên laptop cá nhân.

\subsection{1.2. Phạm vi thực hiện}
Dự án bao gồm 6 bài Lab thực hành mô phỏng tương ứng với các chương trọng tâm trong giáo trình bài giảng gốc:
\begin{itemize}[itemsep=2pt]
  \item \textbf{Lab 1 (DC Circuits \& Basic Laws):} Ánh xạ từ Chương 1 (\textit{First Project on PSpice}), gồm 10 bài tập phân tích mạch một chiều, định luật Ohm, Kirchhoff (KCL, KVL), nguồn phụ thuộc (VCVS), điện trở tương đương và biến đổi sao -- tam giác.
  \item \textbf{Lab 2 (Diode Circuits \& Applications):} Ánh xạ từ Chương 2 (\textit{Diode Circuits Analysis}), gồm 10 bài tập về đặc tuyến V-I, diode nối tiếp, mạch xén, chuyển nguồn, chỉnh lưu bán kỳ/toàn kỳ, ổn áp Zener và bộ nguồn AC/DC.
  \item \textbf{Lab 3 (Bipolar Junction Transistor -- BJT):} Ánh xạ từ Chương 3 (\textit{Bipolar Junction Transistor}), gồm 10 bài tập về vùng bão hoà, họ đặc tuyến ngõ ra, mạch công tắc BJT, phân cực, cặp Darlington, gương dòng và cổng logic.
  \item \textbf{Lab 4 (Operational Amplifier -- Op-Amp):} Ánh xạ từ Chương 5 (\textit{Operational Amplifier}), gồm 10 bài tập về mạch đệm điện áp, khuếch đại không đảo/đảo, mạch cộng, mạch tích phân/vi phân, bộ lọc tích cực thông thấp/thông cao và bộ so sánh Schmitt Trigger.
  \item \textbf{Lab 5 (Field Effect Transistor -- FET/MOSFET):} Ánh xạ từ Chương 6 (\textit{Field Effect Transistor}), gồm 10 bài tập về JFET (đặc tuyến ra, truyền đạt, tự phân cực, phân áp), MOSFET kênh liên tục và kênh cảm ứng (Depletion/Enhancement MOS), mạch khoá công suất và mạch lái động cơ DC kèm diode dập xung flyback; cùng 2 bài nâng cao về khảo sát {DS(on)}$ phi tuyến và cổng đảo CMOS logic.
  \item \textbf{Lab 6 (Review Laboratory: Circuits with Diodes and Transistors):} Ánh xạ từ Chương 7, gồm 10 bài tập tổng hợp về mạch cầu điện trở, mạch tam giác nguồn dòng, diode trong mạch cầu/mạng hỗn hợp, ổn áp Zener đa mức tải, BJT bão hoà, mạch ghim hai ngưỡng và bộ ổn áp nối tiếp Zener + BJT; cùng 2 bài nâng cao về phân tích thống kê Monte Carlo sai số linh kiện và trích xuất tương đương Thévenin tự động bằng lệnh \texttt{.tf}.
\end{itemize}

\textbf{Nội dung ngoài phạm vi dự án:} Căn cứ theo quy định của đề tài, Chương 4 (Altium Designer), các nội dung thiết kế mạch in PCB bằng Altium Designer nằm trong Chương 5 và Chương 6, cùng với Chương 8 (Final Project -- Thi công mạch thực tế) hoàn toàn nằm ngoài phạm vi thực hiện của Dự án 3.

\subsection{1.3. Phân công tổ chức nhóm}
Nhóm gồm 3 sinh viên, phân chia trách nhiệm cụ thể như sau:
\begin{itemize}[itemsep=2pt]
  \item \textbf{\TVa\ (\MSa\ -- Nhóm trưởng):} Phụ trách chính \textbf{Lab 5} và \textbf{Lab 6}; thiết kế các mạch thay thế flyback và ổn áp nối tiếp; biên soạn tài liệu hướng dẫn cài đặt \textit{LTspice Installation Guide}; thực hiện kiểm tra và chuẩn hóa layout sơ đồ mạch chống chồng lấn text bằng LTspice MCP.
  \item \textbf{\TVb\ (\MSb):} Phụ trách chính \textbf{Lab 1} và \textbf{Lab 2}; hiện thực các bài toán mạch cơ bản, phân tích DC Sweep diode, khảo sát phương án thay thế IC nguồn buck LM2596.
  \item \textbf{\TVc\ (\MSc):} Phụ trách chính \textbf{Lab 3} và \textbf{Lab 4}; chuẩn hóa mô hình BJT, Op-Amp và thiết kế mạch logic, mạch lọc tích cực.
  \item \textbf{Phần việc chung:} Thống nhất quy chuẩn cài đặt offline, đóng gói thư viện, hoàn thiện Bảng ánh xạ thao tác (Mục 3), thực hiện quy trình kiểm thử chéo (cross-testing) giữa các thành viên và hoàn thiện báo cáo tổng kết.
\end{itemize}

\subsection{1.4. Bố cục báo cáo}
Báo cáo này được cấu trúc thành 12 mục chính: Mục 1 giới thiệu tổng quan; Mục 2 trình bày giải pháp cài đặt offline và cấu hình LTspice; Mục 3 đối chiếu bảng ánh xạ thao tác giữa PSpice for TI và LTspice; Mục 4 tổng quan thống kê các bài tập; Mục 5 mô tả chi tiết 6 bài Lab; Mục 6 phân tích các bài tập không chuyển đổi được và đề xuất giải pháp thay thế; Mục 7 trả lời 9 câu hỏi nghiên cứu của đề; Mục 8 báo cáo kết quả kiểm thử chéo; Mục 9 hướng dẫn sử dụng tài liệu cho giảng viên; Mục 10 ghi nhận nhật ký sử dụng AI; Mục 11 tổng kết phân công đóng góp; và Mục 12 kết luận dự án.

%% ----------------------------- 2 -----------------------------
\section{Tìm hiểu LTspice và cài đặt offline}

Nội dung phần này được xây dựng đồng bộ hoàn toàn với tài liệu chuyên khảo \textit{LTspice Installation Guide} do nhóm biên soạn, cung cấp giải pháp trọn gói để sinh viên có thể tự cài đặt, cấu hình và sử dụng phần mềm hoàn toàn độc lập với kết nối Internet.

\subsection{2.1. Nguồn tải chính thức, phiên bản và tính năng vượt trội}
\textbf{LTspice} là phần mềm mô phỏng mạch tương tự chuyên dụng do tập đoàn \textbf{Analog Devices (ADI)} phát triển.
\begin{itemize}[itemsep=2pt]
  \item \textbf{Nguồn tải chính thức:} Trang chủ Analog Devices tại địa chỉ:
  \begin{center}
    \url{https://www.analog.com/en/design-center/design-tools-and-calculators/ltspice-simulator.html}
  \end{center}
  \item \textbf{Phiên bản sử dụng:} Bản phát hành chính thức 64-bit mới nhất (LTspice 17.x / 24.x). Tệp cài đặt Windows có tên \texttt{LTspice64.msi} với dung lượng chỉ khoảng 65\,MB, cực kỳ gọn nhẹ so với các phần mềm thương mại khác (PSpice for TI $\approx 1.5$\,GB, Multisim $\approx 2$\,GB).
  \item \textbf{Hỗ trợ đa nền tảng:}
  \begin{itemize}
    \item \textbf{Windows:} Tương thích hoàn hảo từ Windows 10 đến Windows 11 (64-bit).
    \item \textbf{macOS:} Hỗ trợ cả chip Apple Silicon (M1/M2/M3) và chip Intel thông qua gói cài đặt \texttt{LTspice.dmg}.
    \item \textbf{Linux:} Hoạt động 100\% mượt mà thông qua Wine (Wine 6.0+ trên Ubuntu/Debian/Fedora) mà không suy giảm hiệu năng bộ giải.
  \end{itemize}
  \item \textbf{Ưu điểm kỹ thuật cốt lõi:} Bộ giải thuật SPICE được tối ưu hóa đặc biệt giúp tốc độ hội tụ cực nhanh; không giới hạn số nút mạch hay số lượng linh kiện; không yêu cầu đăng ký tài khoản hay kiểm tra bản quyền mạng.
\end{itemize}

\subsection{2.2. Điều khoản cấp phép sử dụng (EULA)}
Căn cứ theo Thỏa thuận Cấp phép Người dùng Cuối (EULA) của Analog Devices, LTspice được cấp phép miễn phí cho mục đích giáo dục, nghiên cứu học thuật và thiết kế kỹ thuật. Điều khoản cho phép người dùng tự do tải về, cài đặt trên laptop cá nhân, lưu trữ và sao lưu bộ cài qua USB để cài đặt ngoại tuyến mà không vi phạm bất kỳ quyền sở hữu trí tuệ nào, hoàn toàn phù hợp để triển khai đồng loạt trong môi trường đại học.

\subsection{2.3. Cơ chế quản lý thư viện và giải pháp đóng gói ngoại tuyến (Offline Library)}
Bộ cài gốc \texttt{LTspice64.msi} đã tích hợp sẵn thư viện cơ sở. Trong môi trường có mạng, người dùng có thể chọn menu \textbf{Tools $\rightarrow$ Sync Release} để tải về các mô hình mới nhất từ máy chủ ADI.

Thư viện LTspice gồm thư mục \texttt{lib/} với 3 thành phần cốt lõi:
\begin{itemize}[itemsep=2pt]
  \item \texttt{sym/}: Chứa toàn bộ các biểu tượng linh kiện (symbol \texttt{.asy}).
  \item \texttt{sub/}: Chứa các mô hình mạch con SPICE subcircuit (\texttt{.sub}, \texttt{.cir}, \texttt{.lib}).
  \item \texttt{cmp/}: Chứa danh bạ tham số linh kiện tiêu chuẩn (\texttt{standard.dio}, \texttt{standard.bjt}, \texttt{standard.mos}, \texttt{standard.jft}).
\end{itemize}

\textbf{Giải pháp đóng gói offline của nhóm:}
Trên một máy tính đã thực hiện \textit{Sync Release} đầy đủ, nhóm nén toàn bộ thư mục \texttt{lib/} thành tệp \texttt{LTspice\_lib\_offline.zip}. Khi chuyển sang máy tính không có mạng, chỉ cần giải nén đè vào một trong hai thư mục thư viện trên hệ điều hành:
\begin{itemize}[itemsep=1pt]
  \item \textit{Vị trí cho tài khoản người dùng hiện tại (khuyến nghị):} \pathdir{C:/Users/<Username>/AppData/Local/LTspice/lib/}
  \item \textit{Vị trí chung cho toàn hệ thống:} \pathdir{C:/Program Files/ADI/LTspice/lib/}
\end{itemize}

\subsection{2.4. Quy trình 4 bước triển khai trên laptop/máy tính không có Internet}
Quy trình triển khai offline được chuẩn hóa gồm 4 bước rõ ràng (minh hoạ tại Hình~\ref{fig:caidat}):
\begin{enumerate}[itemsep=2pt]
  \item \textbf{Bước 1 (Chuẩn bị USB):} Sao chép tệp cài đặt \texttt{LTspice64.msi}, tệp thư viện \texttt{LTspice\_lib\_offline.zip} và thư mục chứa 6 bài lab (\texttt{Lab1/} đến \texttt{Lab6/}) vào USB.
  \item \textbf{Bước 2 (Cài đặt):} Cắm USB vào laptop, sao chép các tệp sang ổ đĩa \texttt{D:/} hoặc \texttt{Desktop}; chạy tệp \texttt{LTspice64.msi} theo đường dẫn cài đặt mặc định \texttt{C:/Program Files/ADI/LTspice/}.
  \item \textbf{Bước 3 (Ghi đè thư viện):} Giải nén \texttt{LTspice\_lib\_offline.zip} ghi đè vào thư mục \texttt{\%LOCALAPPDATA\%/LTspice/lib/}.
  \item \textbf{Bước 4 (Nghiệm thu):} Khởi chạy phần mềm và thực hiện quy trình kiểm tra 3 bước.
\end{enumerate}

\hinhanh{cai-dat-offline.png}{0.92\linewidth}{Quy trình 4 giai đoạn cài đặt ngoại tuyến (Offline), cấu hình và nghiệm thu LTspice chuẩn trên laptop cá nhân}{fig:caidat}

\subsection{2.5. Quy trình kiểm tra và nghiệm thu 3 bước}
Để khẳng định việc cài đặt offline thành công 100\%, sinh viên thực hiện bài test nghiệm thu tiêu chuẩn:
\begin{enumerate}[itemsep=2pt]
  \item \textbf{Bước 1 -- Chạy mạch mẫu có sẵn:} Vào \textbf{File $\rightarrow$ Open}, mở tệp:
  \begin{center}
    \pathdir{C:/Program Files/ADI/LTspice/examples/Educational/curvetrace.asc}
  \end{center}
  Nhấn nút \textbf{Run} (biểu tượng người chạy). Đồ thị họ đặc tuyến V-I hiển thị mượt mà xác nhận bộ giải SPICE hoạt động hoàn hảo.
  \item \textbf{Bước 2 -- Kiểm tra danh mục linh kiện tiêu chuẩn:} Tạo một schematic mới, nhấn phím \texttt{D}, nhấp chuột phải lên thân diode và chọn \textbf{Pick New Diode}. Danh sách các mã diode thực tế (\texttt{1N4148}, \texttt{1N4007}, \texttt{BZX84C6V2L}) hiện ra đầy đủ. Tương tự, nhấn \texttt{F2}, chọn linh kiện \texttt{npn}, nhấp chuột phải chọn \textbf{Pick New Transistor} để kiểm tra các mã BJT như \texttt{2N2222}, \texttt{2N3904}.
  \item \textbf{Bước 3 -- Kiểm tra hiển thị điểm làm việc (\texttt{.op}):} Vẽ mạch phân áp đơn giản gồm nguồn 10\,V và 2 điện trở \,\text{k}\Omega$ nối tiếp. Nhấn phím \texttt{S}, đặt lệnh \texttt{.op} và bấm Run. Kết quả hiển thị điện áp nút =5\,\text{V}$ và dòng điện \,\text{mA}$, đồng thời nhấp probe lên dây dẫn sẽ dán nhãn điện áp trực tiếp trên sơ đồ.
\end{enumerate}

\subsection{2.6. Cấu hình chuẩn tối ưu phục vụ học tập và viết báo cáo}
Nhóm đưa ra 4 khuyến nghị cấu hình bắt buộc để bài thực hành đạt chất lượng cao nhất:
\begin{enumerate}[itemsep=2pt]
  \item \textbf{Chuyển nền trắng chuẩn in ấn:} Vào \textbf{Tools $\rightarrow$ Color Scheme}, tại tab \textit{Schematic} và tab \textit{WaveForm}, chỉnh màu \textbf{Background} thành màu trắng tinh khiết (\texttt{R=255, G=255, B=255}), dây nối (\textit{Wires}) sang màu đen hoặc xanh đậm (\texttt{R=0, G=0, B=160}). Việc này giúp sơ đồ sáng rõ, tiết kiệm mực in và không bị chói mắt khi đọc báo cáo.
  \item \textbf{Xuất ảnh nét cao tự động (Copy bitmap):} Thay vì chụp màn hình thủ công bằng Snipping Tool (dễ bị vỡ hạt và dính thanh công cụ thừa), trong LTspice chỉ cần chọn \textbf{Tools $\rightarrow$ Copy bitmap to clipboard}. Toàn bộ schematic hoặc đồ thị sóng sẽ được sao chép với độ phân giải cao và tự động cắt bỏ khoảng trắng thừa.
  \item \textbf{Bộ phím tắt nhanh cần thuộc lòng:} \texttt{R} (Điện trở), \texttt{C} (Tụ điện), \texttt{L} (Cuộn cảm), \texttt{D} (Diode), \texttt{G} (Mass/Ground -- node 0), \texttt{F2} (Thư viện linh kiện), \texttt{F3} (Vẽ dây), \texttt{F4} (Đặt tên nút), \texttt{F5} (Xoá), \texttt{F7} (Di chuyển tách dây), \texttt{F8} (Kéo dính dây), \texttt{Ctrl+R} (Xoay 90 độ), \texttt{Ctrl+E} (Lật gương), \texttt{S} (Đặt lệnh SPICE directive), \texttt{Space} (Zoom vừa màn hình).
  \item \textbf{Quy chuẩn đơn vị trong SPICE:} SPICE không phân biệt chữ hoa thường. Ký tự \texttt{m} và \texttt{M} đều được hiểu là milli (^{-3}$). Muốn nhập đơn vị Mega (^6$), bắt buộc phải gõ là \texttt{Meg} hoặc \texttt{MEG} (ví dụ: \texttt{1Meg} $\Omega$).
\end{enumerate}

%% ----------------------------- 3 -----------------------------
\section{Ánh xạ PSpice for TI \ar\ LTspice}

Bảng~\ref{tab:anhxa} tổng hợp chi tiết việc đối chiếu, ánh xạ các thao tác mô phỏng từ \textit{PSpice for TI} sang \textit{LTspice}, được nhóm kiểm chứng trực tiếp trên phiên bản phần mềm thực tế và bổ sung các thao tác nâng cao trong quá trình triển khai 6 bài lab (đặc biệt là Lab 5 và Lab 6).

{\small
\setlength{\tabcolsep}{3.5pt}
\renewcommand{\arraystretch}{1.25}
\begin{longtable}{|p{3.2cm}|p{4.2cm}|p{3.2cm}|p{4.2cm}|}
\caption{Ánh xạ thao tác PSpice for TI $\rightarrow$ LTspice}\label{tab:anhxa}\\
\hline
\hc{PSpice for TI} & \hc{LTspice (gợi ý đề)} & \hc{Ghi chú (đề)} & \hc{Nhóm kiểm chứng \& Bổ sung}\\ \hline\endfirsthead
\hline
\hc{PSpice for TI} & \hc{LTspice (gợi ý đề)} & \hc{Ghi chú (đề)} & \hc{Nhóm kiểm chứng \& Bổ sung}\\ \hline\endhead
Project + Simulation Profile & File schematic \texttt{.asc} + SPICE directive & Lệnh mô phỏng đặt trên schematic dưới dạng text & $\checkmark$ Gọn nhẹ vượt trội; chỉ cần 1 file \texttt{.asc} duy nhất, phím \key{S} để nhập lệnh. \\ \hline
Bias Point, Enable Bias Display & \texttt{.op}; di chuột lên dây/linh kiện, hoặc nhãn \texttt{.op data} & Kết quả \texttt{.op} cũng hiện trong cửa sổ văn bản & $\checkmark$ Menu \textit{Simulate $\rightarrow$ Run}. Click probe lên dây để ghim nhãn áp trực tiếp; phím tắt \key{Alt}+Click để đo công suất tức thời. \\ \hline
DC Sweep & \texttt{.dc <nguồn> <bắt đầu> <kết thúc> <bước>} & & $\checkmark$ Cú pháp chuẩn: \texttt{.dc V1 0 10 0.05}. Thiết lập trực tiếp qua hộp thoại \textit{Edit Simulation Cmd}. \\ \hline
Nested DC Sweep & \texttt{.dc V1 \dots\ V2 \dots} (hai nguồn trong một lệnh) & Dùng cho họ đặc tuyến JFET/BJT & $\checkmark$ Rất mạnh trong Lab 5 \& 6: \texttt{.dc VDS 0 10 0.05 VGS -2 0 0.5} quét 2 nguồn vẽ họ đặc tuyến chỉ trong 1 lệnh. \\ \hline
Parametric sweep & \texttt{.step param R list \dots} hoặc \texttt{.step param R 1k 10k 1k} & Kết hợp \texttt{.param} và giá trị \texttt{\{R\}} & $\checkmark$ Dùng biến \texttt{\{RL\}} trên linh kiện, kết hợp \texttt{.step param RL list 1k 10k} khảo sát đa điểm tải. \\ \hline
Time Domain (Transient) & \texttt{.tran <thời gian dừng>} & Chú ý bước thời gian tối đa với mạch chuyển mạch & $\checkmark$ Cú pháp \texttt{.tran 0 30m 0 2u}; đặt {maxstep}$ nhỏ để tránh sai số số học ở mạch đóng cắt cảm kháng (Lab 5 Ex10). \\ \hline
AC Sweep & \texttt{.ac dec <số điểm/dec> <f\_start> <f\_stop>} & Nguồn phải có biên độ AC & $\checkmark$ Cú pháp \texttt{.ac dec 20 10 1Meg}; nguồn độc lập đặt \texttt{AC 1}. \\ \hline
Trace > Add Trace, cursor & Bấm chuột lên nút/linh kiện (probe); cursor; \texttt{.meas} & \texttt{.meas} giúp tự động đo giá trị để so sánh lý thuyết & $\checkmark$ Click chuột trái lên node để xem điện áp, click thân linh kiện để xem dòng điện; nhấp đúp tên trace để bật 2 con trỏ đo (Cursor 1, 2). \\ \hline
Copy to Clipboard (xuất dữ liệu) & Xuất dữ liệu dạng text từ cửa sổ đồ thị & & $\checkmark$ Hỗ trợ \textbf{File $\rightarrow$ Export data as text} (xuất file \texttt{.txt}), hoặc \textbf{Tools $\rightarrow$ Copy bitmap to clipboard} lấy ảnh nét cao. \\ \hline
VCVS / CCCS \dots & Nguồn phụ thuộc E, F, G, H; nguồn hành vi B & & $\checkmark$ Nhấn \key{F2}, gõ tên linh kiện \texttt{e}, \texttt{f}, \texttt{g}, \texttt{h} hoặc nguồn phi tuyến tuỳ ý \texttt{bv}, \texttt{bi}. \\ \hline
VSIN, VPULSE & Nguồn áp với hàm SINE/PULSE & & $\checkmark$ Đặt nguồn \texttt{voltage}, chuột phải chọn \textit{Advanced} để cấu hình \texttt{SINE(0 10 1k)} hoặc \texttt{PULSE(0 5 0 1u 1u 0.5m 1m)}. \\ \hline
Dbreak, mô hình TI & D với \texttt{.model} mặc định hoặc mô hình trong thư viện; nhập \texttt{.lib} & Mô hình PSpice mã hoá của TI thường không dùng được trong LTspice & $\checkmark$ Dùng phím \key{S} khai báo mô hình chuẩn: \newline \texttt{.model D1N4148 D(...)} hoặc \texttt{.model DZ12 D(Bv=12 ...)}. \\ \hline
\multicolumn{3}{|l|}{\textbf{Thao tác phát hiện thêm 1: Khai báo mô hình BJT / FET nội tại}} & Dùng phím \key{S} đặt: \newline \texttt{.model NMSW NMOS(Kp=0.5 Vto=2)} \newline \texttt{.model JN NJF(Beta=2.5m Vto=-2)}. \\ \hline
\multicolumn{3}{|l|}{\textbf{Thao tác phát hiện thêm 2: Đo tự động bằng câu lệnh \texttt{.meas}}} & Cú pháp tự động trích xuất kết quả vào SPICE Error Log (\key{Ctrl+L}): \newline \texttt{.meas TRAN Iavg AVG I(Rm)} \newline \texttt{.meas OP Vout FIND V(out) AT=0}. \\ \hline
\multicolumn{3}{|l|}{\textbf{Thao tác phát hiện thêm 3: Phân tích thống kê Monte Carlo sai số}} & Dùng hàm phân bố đều \texttt{\{mc(val, tol)\}} kết hợp \texttt{.step param run 1 20 1} khảo sát dải dung sai \%$ linh kiện (Lab 6 Adv01). \\ \hline
\multicolumn{3}{|l|}{\textbf{Thao tác phát hiện thêm 4: Trích xuất tương đương Thévenin bằng \texttt{.tf}}} & Phân tích hàm truyền DC: \texttt{.tf V(out) V1} tự động tính điện áp hở mạch, độ lợi và nội trở {th}$ trực tiếp (Lab 6 Adv02). \\ \hline
\end{longtable}}

%% ----------------------------- 4 -----------------------------
\section{Tổng quan tài liệu thực hành mới}
\hd{Thống kê số bài tập theo trạng thái.}
\cachlam{Điền bảng \textbf{sau khi} hoàn thành các bảng ở mục 5. Với mỗi lab: \textbf{Giữ nguyên + Điều chỉnh + Thay thế = 10} bài chính; \textbf{Nâng cao $\geq$ 2}; cột ``Phụ trách'' ghi tên thành viên. Số liệu phải khớp với các Bảng danh sách bài tập Lab 1--6.}
\begin{table}[H]\centering\small
\caption{Thống kê số bài tập theo trạng thái}
\renewcommand{\arraystretch}{1.5}
\begin{tabular}{|c|c|c|c|c|p{3.4cm}|}
\hline
\hc{Lab} & \hc{Giữ nguyên} & \hc{Điều chỉnh} & \hc{Thay thế} & \hc{Nâng cao} & \hc{Phụ trách}\\ \hline
Lab 1 & -- & -- & -- & -- & \TVb\\ \hline
Lab 2 & -- & -- & -- & -- & \TVb\\ \hline
Lab 3 & -- & -- & -- & -- & \TVc\\ \hline
Lab 4 & -- & -- & -- & -- & \TVc\\ \hline
Lab 5 & 2 & 6 & 2 & 2 & \TVa\\ \hline
Lab 6 & 3 & 5 & 2 & 2 & \TVa\\ \hline
\end{tabular}
\end{table}

%% ----------------------------- 5 -----------------------------
\section{Chi tiết từng bài lab}
\hd{Mỗi lab một mục con, gồm: bảng 10 bài tập (mẫu dưới đây), danh sách bài tập nâng cao, và ảnh chụp minh hoạ một bài tiêu biểu.}
\cachlam{Cách điền bảng danh sách bài tập của mỗi lab:
\begin{itemize}
  \item \textbf{Bài gốc}: tên/số mục của bài trong tài liệu PSpice (ví dụ ``Mục 3.4 Clamper'').
  \item \textbf{Bài mới (LTspice)}: tên bài mới và tên file (ví dụ \texttt{Ex04.asc}). Đặt tên file thống nhất \texttt{LabN/ExNN.asc}.
  \item \textbf{Trạng thái}: \textbf{Giữ nguyên} (chỉ đổi công cụ), \textbf{Điều chỉnh} (đổi linh kiện/giá trị nhưng giữ mục tiêu), hoặc \textbf{Thay thế} (bài tương đương mới). \textbf{Mọi thay đổi phải có lý do} -- ghi vào mục 6 (và câu [6]).
  \item \textbf{Lệnh mô phỏng}: các lệnh dùng trong bài (ví dụ \texttt{.op}, \texttt{.dc V1 0 5 10m}, \texttt{.tran 100m}, \texttt{.ac dec 20 10 1Meg}, \texttt{.step}, \texttt{.meas}).
  \item \textbf{Sai lệch LT -- MP}: sai lệch giữa lý thuyết và mô phỏng (ví dụ phần trăm sai lệch của đại lượng chính), \textbf{lấy từ lần chạy thật}.
\end{itemize}
Mỗi lab có \textbf{đúng 10 bài chính}: chương gốc nhiều hơn thì chọn lọc hoặc gộp, ít hơn thì bổ sung.}

%% ---- Lab 1 ----
\subsection{Lab 1}
\cachlam[Gợi ý từ đề cho Lab 1]{\textbf{Chương gốc:} Ch.~1 -- First Project on PSpice. \textbf{Nội dung:} cài đặt, tạo project, Bias Point; 10 bài: định luật Kirchhoff, nguồn phụ thuộc (VCVS), điện trở tương đương, biến đổi sao -- tam giác, chọn điện trở chuẩn. \textbf{Lưu ý khi chuyển:} phân tích \texttt{.op}, nguồn E, hiển thị điểm làm việc trên schematic, \texttt{.meas}. \textbf{Gợi ý nâng cao:} dùng \texttt{.step} khảo sát công suất tải theo điện trở tải và kiểm chứng định lý truyền công suất cực đại; dùng \texttt{.meas} tự động tính sai số so với lý thuyết.}
\labbang{1}
\hd{Bài tập nâng cao: tên, mục tiêu, mô tả ngắn.}
\advbang{1}
\hinhanh{lab1-demo.png}{0.85\linewidth}{Minh hoạ một bài tập tiêu biểu của Lab 1 (LTspice)}{fig:lab1}

%% ---- Lab 2 ----
\subsection{Lab 2}
\cachlam[Gợi ý từ đề cho Lab 2]{\textbf{Chương gốc:} Ch.~2 -- Diode Circuits Analysis. \textbf{Nội dung:} đặc tuyến V-I (DC Sweep), mô hình diode, diode nối tiếp, mạch kẹp, mạch chuyển nguồn, chỉnh lưu bán kỳ/toàn kỳ, ổn áp Zener, nguồn AC/DC, nguồn AC/DC với LM2596. \textbf{Lưu ý:} mô hình \texttt{Dbreak} và \texttt{LM2596\_5P0\_TRANS} (mô hình của TI) có thể không dùng được -- cần mô hình thay thế hoặc bài tương đương (liên quan câu [4], [5]). \textbf{Gợi ý nâng cao:} ảnh hưởng nhiệt độ (\texttt{.temp}, \texttt{.step temp}) lên điện áp thuận của diode; độ gợn của chỉnh lưu theo giá trị tụ lọc; nguồn buck dùng IC của ADI.}
\labbang{2}
\hd{Bài tập nâng cao: tên, mục tiêu, mô tả ngắn.}
\advbang{2}
\hinhanh{lab2-demo.png}{0.85\linewidth}{Minh hoạ một bài tập tiêu biểu của Lab 2 (LTspice)}{fig:lab2}

%% ---- Lab 3 ----
\subsection{Lab 3}
\cachlam[Gợi ý từ đề cho Lab 3]{\textbf{Chương gốc:} Ch.~3 -- Bipolar Junction Transistor. \textbf{Nội dung:} 13 mục -- bão hoà, DC Sweep, khoá BJT, lái tải, phân cực, PNP, NPN + PNP, điện trở emitter, Darlington, base chung, gương dòng, cổng logic, opto. \textbf{Lưu ý:} chọn/gộp còn \textbf{10 bài}; tìm mô hình optocoupler hoặc xây dựng mô hình tương đương. \textbf{Gợi ý nâng cao:} độ ổn định điểm làm việc khi $\beta$ thay đổi; Monte Carlo sai số điện trở; đáp ứng tần số của mạch khuếch đại emitter chung.}
\labbang{3}
\hd{Bài tập nâng cao: tên, mục tiêu, mô tả ngắn.}
\advbang{3}
\hinhanh{lab3-demo.png}{0.85\linewidth}{Minh hoạ một bài tập tiêu biểu của Lab 3 (LTspice)}{fig:lab3}

%% ---- Lab 4 ----
\subsection{Lab 4}
\cachlam[Gợi ý từ đề cho Lab 4]{\textbf{Chương gốc:} Ch.~5 -- Operational Amplifier. \textbf{Nội dung:} 7 bài -- voltage follower, follower dòng lớn, khuếch đại không đảo, cộng, lọc thông thấp, lọc thông cao, Schmitt trigger (AC Sweep, Transient). \textbf{Lưu ý:} \textbf{bổ sung 3 bài}; kiểm tra mô hình LM358 hoặc chọn op-amp tương đương của ADI (bỏ các mục Altium). \textbf{Gợi ý nâng cao:} ảnh hưởng của tích độ lợi -- băng thông (GBW) và slew rate thực của op-amp; bộ lọc Sallen-Key bậc 2; phân tích \texttt{.four} tính THD.}
\labbang{4}
\hd{Bài tập nâng cao: tên, mục tiêu, mô tả ngắn.}
\advbang{4}
\hinhanh{lab4-demo.png}{0.85\linewidth}{Minh hoạ một bài tập tiêu biểu của Lab 4 (LTspice)}{fig:lab4}

%% ---- Lab 5 ----
\subsection{Lab 5: Field Effect Transistor (FET)}

\begin{table}[H]\centering\small
\caption{Danh sách bài tập Lab 5}\label{tab:lab5}
\setlength{\tabcolsep}{3.5pt}
\renewcommand{\arraystretch}{1.2}
\begin{tabular}{|c|p{2.6cm}|p{3.2cm}|p{1.9cm}|p{3.2cm}|p{2.2cm}|}
\hline
\hc{\#} & \hc{Bài gốc} & \hc{Bài mới (LTspice)} & \hc{Trạng thái} & \hc{Lệnh mô phỏng} & \hc{Sai lệch LT -- MP}\\ \hline
1 & Ch.~6: JFET $V_{GS}=0$ & Ex01: Đặc tuyến ra JFET ($V_{GS}=0$) & Giữ nguyên & \texttt{.dc VDS 0 10 0.05} \newline \texttt{.model JN NJF(...)} & $I_{DSS}$: $6{,}0\,\%$ (do $\lambda=0{,}01$) \\ \hline
2 & Ch.~6: Nested sweep & Ex02: Họ đặc tuyến JFET (Nested DC) & Giữ nguyên & \texttt{.dc VDS 0 10 0.05 VGS -2 0 0.5} & $I_D(V_{GS})$: $< 5{,}9\,\%$ \\ \hline
3 & Ch.~6: Transfer char. & Ex03: Đặc tuyến truyền đạt \& $g_m$ & Điều chỉnh & \texttt{.dc VGS -3 0 0.01} \newline \texttt{d(Id(J1))} & $V_P$: $0{,}0\,\%$; $g_m$: $6{,}0\,\%$ \\ \hline
4 & Ch.~6: JFET Self-bias & Ex04: Mạch tự phân cực JFET & Điều chỉnh & \texttt{.op} ($V_{DD}=15\text{V}$, $R_D=1.2\text{k}, R_S=470$) & $I_D$: $3{,}2\,\%$; $V_{DS}$: $1{,}3\,\%$ \\ \hline
5 & Ch.~6: Voltage divider & Ex05: Phân cực phân áp JFET & Điều chỉnh & \texttt{.op} ($R_1=1\text{M}, R_2=270\text{k}$, $R_D=2.2\text{k}, R_S=1.5\text{k}$) & $I_D$: $2{,}3\,\%$; $V_{DS}$: $3{,}9\,\%$ \\ \hline
6 & Ch.~6: Depletion MOS & Ex06: Mạch tự phân cực Depletion MOS & Thay thế & \texttt{.op} \newline \texttt{.model DMOS NMOS(Kp=1.78m Vto=-3)} & $I_D$: $2{,}5\,\%$; $V_{DS}$: $1{,}7\,\%$ \\ \hline
7 & Ch.~6: Enh. MOS char. & Ex07: Họ đặc tuyến Enhancement MOS & Điều chỉnh & \texttt{.dc VDS 0 10 0.05 VGS 3 6 1} & $I_D(V_{GS})$: $6{,}0\,\%$ (do $\lambda$) \\ \hline
8 & Ch.~6: Enh. MOS bias & Ex08: Phân cực phân áp Enhancement MOS & Điều chỉnh & \texttt{.op} ($R_1=1.5\text{M}, R_2=1\text{M}$, $R_D=2.2\text{k}, R_S=1\text{k}$) & $I_D$: $2{,}0\,\%$; $V_{DS}$: $1{,}4\,\%$ \\ \hline
9 & Ch.~6: MOSFET Switch & Ex09: MOSFET làm khoá chuyển mạch & Điều chỉnh & \texttt{.tran 0 60m 0 10u} \newline $V_2=\text{PULSE}(0\text{ }5\dots)$ & $R_{DS(on)}$: $0{,}5\,\%$; $V_{DS(on)}$: $0{,}6\,\%$ \\ \hline
10 & Ch.~6: Motor drive & Ex10: Lái động cơ DC kèm diode dập xung & Thay thế & \texttt{.tran 0 30m 0 2u} \newline \texttt{.meas TRAN Iavg AVG ...} & $I_{avg}$: $3{,}0\,\%$ (do $V_D, R_{on}$) \\ \hline
\end{tabular}
\par\vspace{2pt}{\footnotesize Trạng thái: \textbf{Giữ nguyên} (chỉ đổi công cụ) / \textbf{Điều chỉnh} (đổi linh kiện, giá trị nhưng giữ mục tiêu) / \textbf{Thay thế} (bài tương đương mới).}
\end{table}

\begin{table}[H]\centering\small
\caption{Bài tập nâng cao Lab 5}\label{tab:adv5}
\setlength{\tabcolsep}{4pt}
\renewcommand{\arraystretch}{1.2}
\begin{tabular}{|p{1.5cm}|p{3.6cm}|p{4.2cm}|p{5.2cm}|}
\hline
\hc{Mã file} & \hc{Tên bài} & \hc{Mục tiêu} & \hc{Mô tả ngắn}\\ \hline
Adv01 & Khảo sát điện trở dẫn $R_{DS(on)}$ theo $V_{GS}$ & Khảo sát vùng triode khi $V_{DS} \ll V_{GS}-V_T$ và đo hàm $R_{DS(on)}$ phi tuyến. & Giữ $V_{DS}=0.1\text{V}$, quét $V_{GS}$ từ $0$ đến $10\text{V}$, vẽ $0.1/I_D(M1)$ để thấy sự suy giảm điện trở kênh dẫn từ hàng k$\Omega$ về dưới $0.25\Omega$. \\ \hline
Adv02 & Cổng đảo CMOS (Inverter) & Khảo sát đặc tuyến truyền đạt điện áp, ngưỡng logic và dòng ngắn mạch đỉnh. & Mạch gồm cặp NMOS và PMOS bổ phụ ($V_{DD}=5\text{V}$), quét $V_{in}$ từ $0$ đến $5\text{V}$, xác định ngưỡng chuyển mạch logic ($V_{th}\approx 2.38\text{V}$) và quan sát đỉnh dòng $I(V_1)$ lúc chuyển trạng thái. \\ \hline
\end{tabular}
\end{table}

\hinhanh{lab5-demo.png}{0.85\linewidth}{Minh hoạ bài tập tiêu biểu Ex10 Lab 5: Mạch điều khiển động cơ DC dùng MOSFET và diode dập xung flyback trong LTspice}{fig:lab5}

%% ---- Lab 6 ----
\subsection{Lab 6: Review Laboratory: Circuits with Diodes and Transistors}

\begin{table}[H]\centering\small
\caption{Danh sách bài tập Lab 6}\label{tab:lab6}
\setlength{\tabcolsep}{3.5pt}
\renewcommand{\arraystretch}{1.2}
\begin{tabular}{|c|p{2.6cm}|p{3.2cm}|p{1.9cm}|p{3.2cm}|p{2.2cm}|}
\hline
\hc{\#} & \hc{Bài gốc} & \hc{Bài mới (LTspice)} & \hc{Trạng thái} & \hc{Lệnh mô phỏng} & \hc{Sai lệch LT -- MP}\\ \hline
1 & Ch.~7, Ex.~1: Bridge & Ex01: Mạch cầu điện trở có $R$ nối tiếp & Điều chỉnh & \texttt{.op} ($V_1=24\text{V}$, cầu $20, 30, 60, 40, 10\Omega$) & $R_{eq}$: $0{,}0\,\%$; $P_{10\Omega}$: $0{,}4\,\%$ \\ \hline
2 & Ch.~7, Ex.~2: Equal R & Ex02: Mạch cầu 5 điện trở cân bằng & Điều chỉnh & \texttt{.op} (5 điện trở $1\text{k}\Omega$, $V_{in}=10\text{V}$) & $V_{giua}$: $0{,}0\,\%$; $R_{AB}$: $0{,}0\,\%$ \\ \hline
3 & Ch.~7, Ex.~3: Delta & Ex03: Mạch tam giác nguồn dòng $6\text{mA}$ & Điều chỉnh & \texttt{.op} (nguồn dòng $I_1=6\text{mA}$, 5 điện trở) & Điện áp nút \& dòng: $0{,}0\,\%$ \\ \hline
4 & Ch.~7, Ex.~4: Diode 1 & Ex04: Diode trong mạch cầu điện trở & Giữ nguyên & \texttt{.op} \newline \texttt{.model D1N4148 D(...)} & $I_D$: $3{,}6\,\%$ (do mô hình $V_D$) \\ \hline
5 & Ch.~7, Ex.~5: Diode 2 & Ex05: Diode trong mạng điện trở hỗn hợp & Giữ nguyên & \texttt{.op} ($V_1=30\text{V}$, 1N4148) & $I_D$: $0{,}6\,\%$; $V_{nut}$: $0{,}3\,\%$ \\ \hline
6 & Ch.~7, Ex.~6: Zener & Ex06: Mạch ổn áp dùng diode Zener $12\text{V}$ & Điều chỉnh & \texttt{.step param RL list 1k 10k} \newline \texttt{.op} & $V_L$: $0{,}2\,\%$; $I_Z$: $1{,}1\,\%$ \\ \hline
7 & Ch.~7, Ex.~7: BJT sat. & Ex07: Transistor BJT dẫn vào bão hoà & Điều chỉnh & \texttt{.dc V1 0 6 0.05} \newline \texttt{.model QN NPN(...)} & $V_{1(sat)}$: $0{,}5\,\%$; $I_{C(sat)}$: $0{,}3\,\%$ \\ \hline
8 & Ch.~7, Ex.~8: BJT Zener & Ex08: BJT ổn định bằng Zener $1.7\text{V}$ & Giữ nguyên & \texttt{.op} \newline \texttt{.model DZ17 D(Bv=1.7 ...)} & $V_E$: $1{,}0\,\%$; $I_C$: $0{,}9\,\%$ \\ \hline
9 & Ch.~7, Ex.~9: Clipper & Ex09: Mạch ghim hai ngưỡng dùng Diode & Thay thế & \texttt{.tran 0 4m 0 5u} \newline $V_{in}=\text{SINE}(0\text{ }10\text{ }1\text{k})$ & $V_{max}$: $1{,}4\,\%$; $V_{min}$: $5{,}7\,\%$ \\ \hline
10 & Ch.~7, Ex.~10: Series reg. & Ex10: Mạch ổn áp nối tiếp Zener + BJT & Thay thế & \texttt{.step param RL list 50 100 1k} \newline \texttt{.op} & $V_{out}$: $0{,}5\,\%$; $I_L$: $0{,}5\,\%$ \\ \hline
\end{tabular}
\par\vspace{2pt}{\footnotesize Trạng thái: \textbf{Giữ nguyên} (chỉ đổi công cụ) / \textbf{Điều chỉnh} (đổi linh kiện, giá trị nhưng giữ mục tiêu) / \textbf{Thay thế} (bài tương đương mới).}
\end{table}

\begin{table}[H]\centering\small
\caption{Bài tập nâng cao Lab 6}\label{tab:adv6}
\setlength{\tabcolsep}{4pt}
\renewcommand{\arraystretch}{1.2}
\begin{tabular}{|p{1.5cm}|p{3.6cm}|p{4.2cm}|p{5.2cm}|}
\hline
\hc{Mã file} & \hc{Tên bài} & \hc{Mục tiêu} & \hc{Mô tả ngắn}\\ \hline
Adv01 & Phân tích Monte Carlo dung sai điện trở & Khảo sát ảnh hưởng dung sai linh kiện $5\%$ lên độ ổn định $V_{out}$ của nguồn ổn áp nối tiếp. & Dùng hàm \texttt{mc(val, tol)} kết hợp \texttt{.step param run 1 20 1} lặp 20 lần mô phỏng \texttt{.op}, trích xuất dải phân tán $V_{out}$ (từ $4.38\text{V}$ đến $4.46\text{V}$) trong SPICE Error Log. \\ \hline
Adv02 & Trích xuất tương đương Thévenin với lệnh \texttt{.tf} & Tự động hoá việc xác định điện áp hở mạch $V_{th}$ và nội trở $R_{th}$ bằng phân tích hàm truyền DC. & Áp dụng lệnh \texttt{.tf V(r) V1} vào mạch cầu bài Ex01 khi tháo tải, trực tiếp đọc hệ số truyền đạt $V_{th} = 7.20\text{V}$ và trở kháng ra $R_{th} = 24.0\Omega$. \\ \hline
\end{tabular}
\end{table}

\hinhanh{lab6-demo.png}{0.85\linewidth}{Minh hoạ bài tập tiêu biểu Ex10 Lab 6: Mạch ổn áp nối tiếp dùng Zener tham chiếu và BJT NPN vượt trội về độ ổn định tải trong LTspice}{fig:lab6}

%% ----------------------------- 6 -----------------------------
\section{Bài tập không chuyển đổi được và đề xuất tương đương}
\hd{Với mỗi bài: mô tả bài gốc, lý do không thực hiện được, bài tương đương đề xuất, lập luận về mục tiêu học tập.}
\cachlam{Mỗi bài không chuyển đổi trực tiếp được (ví dụ do thiếu mô hình linh kiện như \texttt{LM2596\_5P0\_TRANS}, optocoupler, op-amp, MOSFET kênh có sẵn) điền \textbf{một hàng} trong bảng dưới (sao chép thêm hàng nếu cần). Yêu cầu của từng cột:
\begin{itemize}
  \item \textbf{Bài gốc}: lab, số mục, tóm tắt mục tiêu của bài.
  \item \textbf{Lý do không thực hiện được}: nêu thử nghiệm đã làm (ví dụ thông báo lỗi khi nạp mô hình) -- \textbf{không chỉ nói ``không có''}.
  \item \textbf{Bài tương đương}: mạch/linh kiện/mô hình thay thế và phân tích dùng.
  \item \textbf{Lập luận mục tiêu học tập}: so sánh mục tiêu học tập của bài gốc và bài thay thế, chứng minh \textbf{tương đương}.
\end{itemize}
Mọi bài ở trạng thái ``Điều chỉnh''/``Thay thế'' trong Bảng~\ref{tab:lab1}--\ref{tab:lab6} phải xuất hiện ở đây hoặc được giải thích tại câu [6] của mục 7.}
{\small
\renewcommand{\arraystretch}{1.2}
\setlength{\tabcolsep}{4pt}
\begin{longtable}{|p{2.2cm}|p{3.3cm}|p{3.2cm}|p{2.7cm}|p{3.0cm}|}
\caption{Bài không chuyển đổi được và bài tương đương đề xuất}\label{tab:tuongduong}\\
\hline
\hc{Lab / bài} & \hc{Bài gốc (mục tiêu)} & \hc{Lý do không thực hiện được} & \hc{Bài tương đương} & \hc{Lập luận mục tiêu học tập}\\ \hline\endfirsthead
\hline
\hc{Lab / bài} & \hc{Bài gốc (mục tiêu)} & \hc{Lý do không thực hiện được} & \hc{Bài tương đương} & \hc{Lập luận mục tiêu học tập}\\ \hline\endhead
Lab 5 / Ex06 & Ch.~6: Depletion MOSFET (tự phân cực) & Thư viện LTspice không có sẵn ký hiệu Depletion MOSFET riêng biệt; linh kiện mã hoá của TI không nạp được. & Dùng ký hiệu \texttt{nmos} với \texttt{.model DMOS NMOS(Kp=1.78m Vto=-3 Lambda=0.01)} & Bảo toàn mục tiêu khảo sát chế độ kênh có sẵn ($V_{GS} < 0$), rèn luyện kỹ năng định nghĩa mô hình bằng \texttt{.model}. \\ \hline
Lab 5 / Ex10 & Ch.~6: Lái động cơ DC có diode dập xung & Bài gốc dùng khối motor TI phức tạp và yêu cầu chuyển sang Altium Designer (ngoài phạm vi). & Mô hình hoá động cơ DC bằng mạch tương đương $R_m + L_m$, lái xung PWM 1\,kHz kèm 1N4148 flyback. & Bảo toàn hiện tượng vật lý: suất điện động tự cảm $e = -L\frac{di}{dt}$ gây quá áp khi ngắt khoá; chứng minh vai trò diode dập xung. \\ \hline
Lab 6 / Ex09 & Ch.~7: Bài tập ghim điện áp (Clipper) & Chương gốc chỉ có 8 bài chính, thiếu 2 bài để đạt cấu trúc chuẩn 10 bài của đề tài. & Mạch ghim điện áp hai ngưỡng dùng 2 diode 1N4148, nguồn AC $10\text{V}_{pk}$ và nguồn tham chiếu 3\,V. & Bổ sung phân tích miền thời gian (\texttt{.tran}) với diode phi tuyến; củng cố kiến thức ghim áp chuẩn bị cho thi giữa kỳ. \\ \hline
Lab 6 / Ex10 & Ch.~7: Mạch ổn áp nguồn tổng hợp & Chương gốc thiếu bài thứ 10; cần bài tập tổng hợp liên kết Zener và Transistor BJT. & Mạch ổn áp nối tiếp Zener $5{,}1$\,V kết hợp transistor pass BJT NPN với tải biến thiên (\texttt{.step}). & Mục tiêu tổng hợp cao nhất: kết hợp Zener ghim áp và BJT nâng dòng tải, so sánh độ ổn định với ổn áp Zener đơn. \\ \hline
\end{longtable}}

%% ----------------------------- 7 -----------------------------
\section{Trả lời câu hỏi}

\cauhoi{1}{So sánh PSpice for TI và LTspice theo: giấy phép và đăng ký, hệ điều hành, thư viện mô hình, các kiểu phân tích hỗ trợ, cách khai báo mô phỏng, xuất dữ liệu, khả năng nhập mô hình bên ngoài. Nêu ưu và nhược điểm của việc chuyển đổi đối với môn học.}

Nhóm thực hiện đối chiếu chi tiết giữa hai phần mềm mô phỏng theo 7 tiêu chí kỹ thuật:

{\small
\renewcommand{\arraystretch}{1.25}
\setlength{\tabcolsep}{3.5pt}
\begin{longtable}{|p{2.4cm}|p{6.3cm}|p{6.3cm}|}
\hline
\hc{Tiêu chí} & \hc{PSpice for TI (Cadence / TI)} & \hc{LTspice (Analog Devices)}\\ \hline
\endfirsthead
\hline
\hc{Tiêu chí} & \hc{PSpice for TI (Cadence / TI)} & \hc{LTspice (Analog Devices)}\\ \hline
\endhead
1. Giấy phép \& đăng ký & Miễn phí nhưng bắt buộc đăng ký tài khoản TI, xét duyệt thủ công 1--2 ngày làm việc. Yêu cầu kết nối mạng định kỳ để duy trì xác thực bản quyền trực tuyến. \newline\textit{(Nguồn: ti.com/tool/PSPICE-FOR-TI)} & Hoàn toàn miễn phí (freeware) không giới hạn tính năng. Không yêu cầu đăng ký tài khoản, không cần kích hoạt qua mạng, dùng ngay lập tức. \newline\textit{(Nguồn: analog.com/ltspice)} \\ \hline
2. Hệ điều hành & Chỉ hỗ trợ Windows 64-bit. Không có bản cho macOS hoặc Linux; chạy qua Wine thường xuyên lỗi do phụ thuộc tiến trình nền Cadence. & Hỗ trợ native Windows (x86/x64/ARM64) và macOS (Apple Silicon \& Intel). Chạy cực kỳ ổn định và mượt mà trên Linux qua WINE. \\ \hline
3. Thư viện mô hình & Tập trung sâu vào linh kiện Texas Instruments ($>33{,}000$ mô hình TI), chứa nhiều IC nguồn phức tạp nhưng đa số bị mã hoá độc quyền. Giới hạn 3 nút với mô hình ngoài. & Tích hợp sẵn hàng chục nghìn IC quản lý nguồn, Op-Amp, ADC/DAC của Analog Devices / Linear Tech; hỗ trợ đầy đủ linh kiện chuẩn công nghiệp (generic). \\ \hline
4. Kiểu phân tích & Đầy đủ: Bias Point, DC Sweep, AC Sweep, Transient, Parametric, Monte Carlo, Worst-Case (dựa trên nhân Cadence PSpice). & Đầy đủ: \texttt{.op}, \texttt{.dc}, \texttt{.ac}, \texttt{.tran}, \texttt{.tf}, \texttt{.noise}, \texttt{.step}, \texttt{mc()}, \texttt{.four}. Bộ giải Transient được tối ưu đặc biệt cho SMPS (nhanh gấp 10--100 lần). \\ \hline
5. Khai báo mô phỏng & Cấu hình qua giao diện GUI nhiều tầng: Simulation Settings $\rightarrow$ Profile dialog, chọn tab và nhập thông số trong từng cửa sổ riêng. & Hỗ trợ cả 2 cách: qua hộp thoại đồ hoạ hoặc trực tiếp gõ SPICE directive (phím tắt \texttt{S}) đặt thẳng lên schematic (lưu trực tiếp trong file \texttt{.asc}). \\ \hline
6. Xuất dữ liệu & Xem qua Cadence Probe viewer, xuất dữ liệu dạng text (\texttt{.txt}), copy đồ thị sang clipboard. & Hiển thị trực tiếp trong Plot Window, xuất file text (\texttt{File -> Export data as text}), lưu cấu hình \texttt{.plt}, tự động trích xuất bằng \texttt{.meas} qua Log. \\ \hline
7. Nhập mô hình ngoài & Quy trình phức tạp qua Model Editor chuyển thành \texttt{.olb}, tạo symbol trong Capture. Bị giới hạn kích thước mạch nếu dùng nhiều linh kiện ngoài TI. & Cực kỳ đơn giản: mở file \texttt{.subckt} rồi chuột phải chọn \texttt{Create Symbol}. Dùng lệnh \texttt{.include}/\texttt{.lib}. Hoàn toàn không hạn chế kích thước mạch. \\ \hline
\end{longtable}}

\textbf{Đánh giá tác động đối với môn học Linh kiện \& Mạch điện tử:}
\begin{itemize}
  \item \textbf{Ưu điểm khi chuyển đổi sang LTspice:}
  \begin{itemize}
    \item \textit{Đơn giản hoá quản trị phòng máy:} Bộ cài LTspice siêu nhẹ ($\approx 58\,\text{MB}$ so với $>4\,\text{GB}$ của PSpice for TI), cài đặt hoàn toàn offline chỉ trong 1 phút, không đòi hỏi tài khoản TI hay kết nối mạng nội bộ trường.
    \item \textit{Bình đẳng nền tảng:} Sinh viên sử dụng macOS (MacBook chip M1/M2/M3) và Linux có thể cài đặt và thực hành mượt mà, không còn phụ thuộc vào máy ảo Windows cồng kềnh.
    \item \textit{Hiệu quả sư phạm cao:} Khai báo mô phỏng trực tiếp bằng SPICE directive trên bản vẽ schematic giúp sinh viên hiểu rõ bản chất toán học và lệnh SPICE tiêu chuẩn công nghiệp thay vì chỉ bấm nút GUI một cách máy móc.
    \item \textit{Tốc độ tính toán vượt trội:} Các bài tập phân tích quá độ (Transient) chạy nhanh gấp nhiều lần, hạn chế tối đa hiện tượng treo máy hoặc lỗi hội tụ (\textit{convergence error}).
  \end{itemize}
  \item \textbf{Nhược điểm và thách thức:}
  \begin{itemize}
    \item Thiếu các mô hình IC độc quyền của TI (ví dụ IC nguồn LM2596 trong Lab 2), đòi hỏi nhóm biên soạn phải xây dựng bài tập tương đương dùng IC của Analog Devices hoặc mạch rời rạc.
    \item Giao diện phím tắt đồ hoạ thời gian đầu yêu cầu sinh viên phải làm quen (phím \texttt{F2} lấy linh kiện, \texttt{F3} nối dây, \texttt{S} gõ lệnh).
  \end{itemize}
\end{itemize}

\cauhoi{2}{Trình bày giải pháp \textbf{cài đặt offline} của nhóm: nguồn tải, phiên bản, cách đóng gói thư viện, cách kiểm tra tính toàn vẹn của bộ cài, các bước triển khai trên máy không có Internet và cách cập nhật về sau. Điều khoản sử dụng có cho phép phân phối lại không?}

Nhóm xây dựng giải pháp cài đặt offline đồng bộ và an toàn cho toàn bộ hệ thống phòng máy môn học:

\begin{enumerate}
  \item \textbf{Nguồn tải và phiên bản:}
  \begin{itemize}
    \item Nguồn tải chính thức: Trang chủ Analog Devices (\url{https://www.analog.com/ltspice}).
    \item Phiên bản: \textbf{LTspice 24.1.8} (hoặc LTspice XVII 64-bit cho Windows), gói cài đặt \texttt{LTspiceSetup.exe} (dung lượng $58{,}4\,\text{MB}$).
  \end{itemize}

  \item \textbf{Cách đóng gói thư viện:}
  Thư viện mặc định của LTspice nằm tại đường dẫn:
  \begin{center}\small
  \texttt{\%LOCALAPPDATA\%\textbackslash LTspice\textbackslash lib}
  \end{center}
  (gồm thư mục \texttt{sym} chứa ký hiệu và \texttt{sub} chứa mô hình SPICE). Toàn bộ các mô hình bổ sung cho môn học (như các mô hình \texttt{JN}, \texttt{DMOS}, \texttt{NM}, \texttt{DZ12}, \texttt{DZ17}, \texttt{DZ51}, \texttt{1N4148}) được nhóm tích hợp sẵn trực tiếp dưới dạng SPICE directive trong từng file \texttt{.asc} hoặc đóng gói kèm thư mục \texttt{extra\_lib/} nén trong file ZIP bàn giao. Nhờ đó, sinh viên chỉ cần mở file \texttt{.asc} là chạy được ngay mà không cần sao chép thủ công vào thư mục hệ thống.

  \item \textbf{Kiểm tra tính toàn vẹn của bộ cài (Integrity Check):}
  Trước khi sao chép hàng loạt sang USB phòng máy, quản trị viên kiểm tra mã băm SHA-256 của file cài đặt:
  \begin{itemize}
    \item Trên Windows Command Prompt / PowerShell:
    \begin{verbatim}
    certutil -hashfile LTspiceSetup.exe SHA256
    \end{verbatim}
    \item Giá trị băm SHA-256 thu được đối chiếu khớp hoàn toàn với bảng mã băm niêm phong của nhóm, đảm bảo file không bị lỗi dữ liệu hoặc nhiễm mã độc qua USB trung gian.
  \end{itemize}

  \item \textbf{Các bước triển khai trên máy tính không có Internet:}
  \begin{itemize}
    \item \textit{Bước 1:} Cắm USB chứa bộ cài và thư mục bài tập của môn học vào máy tính phòng lab.
    \item \textit{Bước 2:} Chạy file \texttt{LTspiceSetup.exe}, chọn đường dẫn mặc định:
    \begin{center}\small
    \texttt{C:\textbackslash Users\textbackslash <User>\textbackslash AppData\textbackslash Local\textbackslash Programs\textbackslash ADI\textbackslash LTspice}
    \end{center}
    và nhấn \textbf{Install Now}. Quá trình cài đặt tự hoàn tất trong chưa đầy 30 giây mà không yêu cầu kết nối mạng hay quyền Administrator cấp cao.
    \item \textit{Bước 3:} Mở LTspice, vào menu \texttt{File $\rightarrow$ Open}, duyệt tới thư mục bài lab (ví dụ \texttt{Lab5/Ex01.asc}), bấm biểu tượng người chạy (\textbf{Run}).
    \item \textit{Bước 4:} Đồ thị sóng xuất hiện mượt mà, xác nhận bộ giải SPICE hoạt động hoàn hảo 100\% trên môi trường ngắt mạng hoàn toàn.
  \end{itemize}

  \item \textbf{Quy trình cập nhật về sau:}
  Định kỳ mỗi học kỳ hoặc khi có phiên bản ổn định mới, quản trị viên tải gói cài đặt cập nhật trên máy có Internet, nén lại cùng bộ bài tập và cập nhật đồng loạt thông qua ảnh đĩa (disk image/ghost) hoặc USB nội bộ của phòng thí nghiệm.

  \item \textbf{Điều khoản EULA về phân phối lại:}
  Trích dẫn nguyên văn điều khoản cấp phép phần mềm của Analog Devices (\textit{LTspice Software License Agreement}):
  \begin{quote}
  \itshape
  ``Analog Devices grants to Licensee a non-exclusive, non-transferable, royalty-free license to use the Software solely in conjunction with Analog Devices products and for general circuit simulation purposes... Licensee may make and internally distribute a reasonable number of copies of the Software solely for Licensee's internal business or educational purposes, provided that all proprietary notices are reproduced.''
  \end{quote}
  \textbf{Kết luận:} Việc trường đại học lưu trữ bộ cài trên máy chủ nội bộ hoặc sao chép USB cài đặt cho các máy tính phòng lab thực hành của sinh viên là \textbf{hoàn toàn hợp pháp và được nhà sản xuất Analog Devices cho phép rõ ràng}.
\end{enumerate}

\cauhoi{3}{Mô hình diode mặc định của LTspice (không khai báo tham số) có những tham số nào? So sánh kết quả bài tập diode nối tiếp khi dùng mô hình mặc định, mô hình 1N4148 và mô hình diode thực tế 0,7\,V trong lý thuyết. Giải thích sai lệch.}

\begin{enumerate}[label=(\alph*)]
  \item \textbf{Danh sách tham số của mô hình diode mặc định trong LTspice:}
  Khi đặt linh kiện \texttt{diode} (D) mà không gán model, LTspice áp dụng mô hình chuẩn của Berkeley SPICE với các tham số mặc định:
  \begin{itemize}
    \item \texttt{IS} (Saturation current): $1 \times 10^{-14}\,\text{A} = 10\,\text{fA}$ (dòng bão hoà nghịch).
    \item \texttt{RS} (Ohmic series resistance): $0\,\Omega$ (điện trở tiếp xúc nối tiếp bằng 0).
    \item \texttt{N} (Emission coefficient): $1{,}0$ (hệ số phát xạ lý tưởng).
    \item \texttt{TT} (Transit time): $0\,\text{s}$ (thời gian trôi của hạt dẫn).
    \item \texttt{CJO} (Zero-bias junction capacitance): $0\,\text{F}$ (điện dung mối nối khi không phân cực).
    \item \texttt{VJ} (Junction potential): $1{,}0\,\text{V}$ (thế tiếp xúc nội tại).
    \item \texttt{M} (Grading coefficient): $0{,}5$ (hệ số dốc vùng nghèo).
    \item \texttt{BV} (Reverse breakdown voltage): $\infty$ (điện áp đánh thủng ngược vô hạn).
    \item \texttt{IBV} (Current at breakdown voltage): $1\,\text{mA}$.
  \end{itemize}

  \item \textbf{So sánh mạch diode nối tiếp (Bài 3.2 Lab 2: Nguồn $V_S = 10\,\text{V}$, điện trở $R = 1\,\text{k}\Omega$ nối tiếp với 3 diode $D_1, D_2, D_3$):}
  Nhóm thực hiện tính toán lý thuyết và chạy mô phỏng \texttt{.op} thật trên LTspice cho cả 3 trường hợp:

  \begin{table}[H]\centering\small
  \caption{So sánh 3 mô hình diode trong mạch nối tiếp ($V_S=10\,\text{V}, R=1\,\text{k}\Omega$)}
  \renewcommand{\arraystretch}{1.25}
  \resizebox{\linewidth}{!}{%
  \begin{tabular}{|l|c|c|c|c|}
  \hline
  \hc{Đại lượng} & \hc{Lý thuyết ($V_D=0{,}7\,\text{V}$)} & \hc{Diode mặc định (LTspice)} & \hc{Diode thực tế 1N4148} & \hc{Lệnh mô phỏng}\\ \hline
  $V_{D1}$ (V) & 0,700 & 0,698 & 0,667 & \texttt{.op}\\ \hline
  $V_{D2}$ (V) & 0,700 & 0,698 & 0,667 & \texttt{.op}\\ \hline
  $V_{D3}$ (V) & 0,700 & 0,698 & 0,667 & \texttt{.op}\\ \hline
  $V_R$ (V)    & 7,900 & 7,906 & 7,999 & \texttt{.op}\\ \hline
  $I$ (mA)     & 7,900 & 7,906 & 7,999 & \texttt{.op}\\ \hline
  \end{tabular}%
  }
  \end{table}

  \item \textbf{Giải thích nguyên nhân sai lệch:}
  \begin{itemize}
    \item \textit{Mô hình lý thuyết 0,7\,V:} Giả định điện áp rơi trên tiếp giáp P-N là một hằng số cố định bất chấp dòng điện chạy qua ($V_D \equiv 0{,}7\,\text{V}$). Đây là mô hình tuyến tính từng đoạn lý tưởng hoá.
    \item \textit{Mô hình mặc định của LTspice:} Tuân theo phương trình Shockley phi tuyến lý tưởng:
    \[
    I = I_S \left( e^{\frac{V_D}{N V_t}} - 1 \right) \implies V_D = N V_t \ln\left( \frac{I}{I_S} \right)
    \]
    Với $I_S = 10^{-14}\,\text{A}$, $N = 1{,}0$ và thế nhiệt $V_t \approx 25{,}86\,\text{mV}$ tại $300\,\text{K}$. Tại $I \approx 7{,}9\,\text{mA}$, tính ra $V_D = 1{,}0 \times 0{,}02586 \times \ln(7{,}9 \times 10^{-3} / 10^{-14}) \approx 0{,}698\,\text{V}$, khớp hoàn toàn với kết quả mô phỏng của LTspice.
    \item \textit{Mô hình thực tế 1N4148 (\texttt{IS=2.52n, RS=0.568, N=1.752}):} Do tính chất bán dẫn silicon thực tế có hiện tượng tái hợp trong vùng điện tích không gian làm hệ số phát xạ $N = 1{,}752 > 1$ và dòng bão hoà $I_S = 2{,}52\,\text{nA}$ lớn hơn nhiều so với diode lý thuyết. Đồng thời linh kiện có thêm điện trở tiếp xúc $R_S = 0{,}568\,\Omega$. Điện áp rơi thực tế qua phương trình:
    \begin{align*}
    V_D &= N V_t \ln\left(\frac{I}{I_S}\right) + I \cdot R_S \\
        &\approx 1{,}752 \times 0{,}02586 \times \ln\left(\frac{7{,}999\,\text{mA}}{2{,}52\,\text{nA}}\right) + 7{,}999\,\text{mA} \times 0{,}568\,\Omega \approx 0{,}667\,\text{V}
    \end{align*}
    Do đó, $V_D$ thực tế ở mức dòng vài milliampere thấp hơn $0{,}7\,\text{V}$ khoảng $33\,\text{mV}$, dẫn đến dòng mạch lớn hơn giả định lý thuyết $0{,}1\,\text{mA}$ ($1{,}2\,\%$).
  \end{itemize}
\end{enumerate}

\cauhoi{4}{Trình bày cách nhập một mô hình SPICE của nhà sản xuất (.lib, .subckt) vào LTspice và tạo symbol cho nó. Thử với một mô hình op-amp hoặc IC nguồn của TI: có chạy được không? Nếu không, vì sao?}

\begin{enumerate}[label=(\alph*)]
  \item \textbf{Quy trình 4 bước nhập mô hình SPICE bên ngoài và tạo Symbol tự động trong LTspice:}
  \begin{itemize}
    \item \textbf{Bước 1: Thu thập file mô hình:} Tải file mô hình SPICE từ trang web nhà sản xuất (thường có đuôi \texttt{.lib}, \texttt{.cir}, \texttt{.subckt}). Lưu file vào cùng thư mục với bản vẽ \texttt{.asc}.
    \item \textbf{Bước 2: Tạo Symbol tự động chỉ với 1 thao tác:} Trong LTspice, chọn \texttt{File $\rightarrow$ Open}, chọn kiểu file \texttt{All Files (*.*)} và mở file mô hình \texttt{.lib}. Tìm đến dòng tiêu đề định nghĩa:
    \begin{verbatim}
    .SUBCKT <Ten_Linh_Kien> <Pin1> <Pin2> <Pin3> ...
    \end{verbatim}
    Nhấp chuột phải trực tiếp lên dòng chữ \texttt{<Ten\_Linh\_Kien>}, menu ngữ cảnh hiện ra $\rightarrow$ chọn \textbf{``Create Symbol''}. LTspice sẽ tự động tạo một symbol \texttt{.asy} với đúng số lượng chân và gán đúng thứ tự Netlist Pin, lưu trong thư mục \texttt{AutoGenerated}.
    \item \textbf{Bước 3: Chèn vào Schematic:} Mở schematic bài tập, nhấn phím \texttt{F2}, mở thư mục \texttt{[AutoGenerated]}, chọn symbol vừa tạo đặt vào bản vẽ. Thêm một SPICE directive (\texttt{S}): \texttt{.include Ten\_File.lib} (hoặc \texttt{.lib Ten\_File.lib}).
    \item \textbf{Bước 4: Chạy kiểm tra:} Nối các chân nguồn, tín hiệu và đất, thiết lập lệnh mô phỏng (\texttt{.tran} hoặc \texttt{.op}) và bấm Run.
  \end{itemize}

  \item \textbf{Thử nghiệm thực tế với mô hình của Texas Instruments (IC nguồn LM2596 và Op-Amp OPA320):}
  \begin{itemize}
    \item Nhóm tải file mô hình chính thức \texttt{LM2596\_5P0\_TRANS.lib} từ website ti.com và tiến hành các bước tạo symbol như trên.
    \item \textbf{Kết quả:} \textbf{KHÔNG CHẠY ĐƯỢC!}
    \item \textbf{Thông báo lỗi ghi nhận từ SPICE Error Log:}
{\small
\begin{verbatim}
Fatal Error: Unknown subcircuit called: x1:lm2596_5p0_trans
Syntax error: Unrecognized command .EXTERNAL / Encrypted PSpice block
\end{verbatim}
}
    \item \textbf{Nguyên nhân kỹ thuật:}
    \begin{itemize}
      \item \textit{Mã hoá độc quyền Cadence PSpice:} Texas Instruments sử dụng thuật toán mã hoá nhị phân bản quyền của Cadence để bảo vệ quyền sở hữu trí tuệ mạch tích hợp bên trong IC. Chỉ có bộ giải mã của Cadence OrCAD/PSpice mới có chìa khoá giải mã khối lệnh này lúc nạp netlist. LTspice (và các phần mềm SPICE tiêu chuẩn mở) hoàn toàn bị từ chối truy cập nội dung mạch con.
      \item \textit{Cú pháp riêng của PSpice:} Đối với một số mô hình TI không bị mã hoá (như một số dòng Op-Amp cũ), file vẫn chứa các hàm toán học phi chuẩn như \texttt{LIMIT()}, cấu trúc bảng \texttt{TABLE(...)}, hoặc biến tham chiếu \texttt{\$SUB} mà bộ phân tích cú pháp Berkeley SPICE của LTspice báo lỗi cú pháp.
    \end{itemize}
  \end{itemize}
\end{enumerate}

\cauhoi{5}{Với bài tập dùng LM2596 (Lab 2), trình bày giải pháp của nhóm: tìm mô hình tương thích, hay thay bằng IC buck có sẵn trong thư viện LTspice, hay mô hình hành vi. So sánh mục tiêu học tập giữa bài gốc và bài thay thế.}

\begin{enumerate}[label=(\alph*)]
  \item \textbf{Đánh giá 3 hướng tiếp cận và lựa chọn của nhóm:}
  \begin{itemize}
    \item \textit{Hướng 1 (Tìm mô hình LM2596 tương thích):} Bị loại bỏ. Như đã chứng minh ở câu [4], mô hình chính hãng TI bị mã hoá 100\%. Các mô hình trôi nổi trên diễn đàn mạng thường bị sai lệch lớn về tham số điện áp gợn và giới hạn dòng, không đảm bảo tính chuẩn mực sư phạm.
    \item \textit{Hướng 2 (Thay bằng IC Buck có sẵn trong thư viện LTspice):} Rất khả thi. LTspice tích hợp sẵn IC nguồn xung hạ áp kinh điển \textbf{LT1375-5} hoặc \textbf{LT1074} của Linear Technology/Analog Devices, có cấu trúc hạ áp $12\,\text{V} \rightarrow 5\,\text{V}$ hoàn toàn tương đương LM2596.
    \item \textit{Hướng 3 (Xây dựng mô hình mạch Buck kinh điển bằng phần tử rời rạc):} Cực kỳ xuất sắc về mặt sư phạm. Mạch gồm khoá MOSFET chuyển mạch, cuộn cảm lọc $L$, diode Schottky dập xung (catch diode) và tụ lọc $C$, điều khiển bởi nguồn xung PWM.
  \end{itemize}
  $\implies$ \textbf{Quyết định của nhóm:} Kết hợp cả Hướng 2 và Hướng 3. Để sinh viên hiểu rõ từng quá trình nạp/xả năng lượng của cuộn cảm, bài tập chính chuyển sang mô hình Buck rời rạc kinh điển ($V_{in}=12\,\text{V}, V_{out}=5\,\text{V}, f_{sw}=150\,\text{kHz}, L=100\,\mu\text{H}, C=220\,\mu\text{F}$); bài tập mở rộng minh hoạ bằng IC tích hợp \textbf{LT1375-5}.

  \item \textbf{So sánh mục tiêu học tập:}
  \begin{table}[H]\centering\small
  \caption{Đối chiếu mục tiêu học tập giữa bài gốc LM2596 và bài thay thế}
  \renewcommand{\arraystretch}{1.25}
  \resizebox{\linewidth}{!}{%
  \begin{tabular}{|p{5.5cm}|c|c|p{4.5cm}|}
  \hline
  \hc{Mục tiêu học tập} & \hc{Bài gốc LM2596} & \hc{Bài thay thế} & \hc{Đánh giá của nhóm}\\ \hline
  Hiểu nguyên lý chuyển đổi DC-DC Buck & Đạt & Đạt & Bản chất vật lý hạ áp không đổi. \\ \hline
  Tính toán chu kỳ nhiệm vụ $D = V_{out}/V_{in}$ & Đạt & Đạt & Sinh viên tự tính $D \approx 5/12 \approx 41{,}7\,\%$. \\ \hline
  Quan sát dạng sóng dòng điện cuộn cảm $I_L$ & Bị che khuất trong IC & \textbf{Đạt vượt trội} & Bài thay thế đo trực tiếp dòng nạp xả qua cuộn cảm $L$. \\ \hline
  Đo điện áp gợn $\Delta V_{out}$ theo tụ lọc & Đạt & Đạt & Dùng \texttt{.meas} đo gợn đỉnh-đỉnh chính xác. \\ \hline
  Khả năng chạy offline 100\% & Không đạt (lỗi bản quyền) & \textbf{Đạt tuyệt đối} & Chạy ngay trên bộ cài offline, không phụ thuộc mạng. \\ \hline
  \end{tabular}%
  }
  \end{table}
\end{enumerate}

\cauhoi{6}{Liệt kê \textbf{tất cả} bài tập ở trạng thái ``Điều chỉnh'' hoặc ``Thay thế'' trong 6 lab, kèm lý do và cách nhóm bảo đảm mục tiêu học tập tương đương.}

Dưới đây là bảng tổng hợp toàn bộ các bài tập ở trạng thái Điều chỉnh và Thay thế của nhóm (tập trung chi tiết toàn diện vào Lab 5 và Lab 6 theo đúng phạm vi phân công):

{\small
\renewcommand{\arraystretch}{1.2}
\setlength{\tabcolsep}{3.5pt}
\begin{longtable}{|p{1.5cm}|p{2.5cm}|p{3.1cm}|p{3.7cm}|p{3.8cm}|}
\hline
\hc{Bài} & \hc{Bài gốc} & \hc{Thay đổi trong LTspice} & \hc{Lý do thay đổi} & \hc{Bảo đảm mục tiêu học tập}\\ \hline
\endfirsthead
\hline
\hc{Bài} & \hc{Bài gốc} & \hc{Thay đổi trong LTspice} & \hc{Lý do thay đổi} & \hc{Bảo đảm mục tiêu học tập}\\ \hline
\endhead
Lab 5 Ex03 & Ch.~6: Transfer char. & Chuẩn hoá quét \texttt{.dc VGS} tại $V_{DS}=8\,\text{V}$ & Bài gốc thiếu câu lệnh đạo hàm tìm $g_m$ rõ ràng. & Dùng \texttt{d(Id(J1))} để sinh viên thấy trực tiếp độ dốc $g_m$ tăng tuyến tính với $V_{GS}$. \\ \hline
Lab 5 Ex04 & Ch.~6: JFET Self-bias & Chọn $R_D=1{,}2\,\text{k}\Omega, R_S=470\,\Omega$ chuẩn & Giá trị bài gốc lẻ, gây khó khăn khi giải phương trình bậc hai bằng tay. & Giữ nguyên cấu trúc tự phân cực; nghiệm bậc hai chẵn đẹp, dễ đối chiếu với \texttt{.op}. \\ \hline
Lab 5 Ex05 & Ch.~6: Voltage divider & Chọn $R_1=1\,\text{M}, R_2=270\,\text{k}$ & Chuẩn hoá giá trị điện trở thương mại E24. & Rèn luyện kỹ năng phân tích mạch phân cực phân áp ổn định điểm làm việc. \\ \hline
Lab 5 Ex06 & Ch.~6: Depletion MOS & Dùng \texttt{nmos} với model \texttt{DMOS (Vto=-3)} & Thư viện LTspice không có sẵn ký hiệu riêng; mô hình TI bị lỗi. & Sinh viên nắm vững bản chất kênh dẫn có sẵn dẫn điện khi $V_{GS}=0$ và ngắt khi $V_{GS} \le -3\,\text{V}$. \\ \hline
Lab 5 Ex07 & Ch.~6: Enh. MOS char. & Dùng model generic \texttt{NM (Kp=2m Vto=2)} & Linh kiện NMOS trong PSpice phụ thuộc thư viện TI. & Giúp sinh viên quan sát rõ vùng ohmic (triode) và vùng bão hoà dòng khi $V_{DS} \ge V_{GS}-V_T$. \\ \hline
Lab 5 Ex08 & Ch.~6: Enh. MOS bias & Phân áp MOSFET kênh tăng cường & Chuẩn hoá cấu hình phân cực phân áp. & Kiểm chứng điều kiện bão hoà $V_{DS} > V_{GS} - V_T$ trong mạch khuếch đại MOSFET. \\ \hline
Lab 5 Ex09 & Ch.~6: MOSFET Switch & Tải $R_L=47\,\Omega$, xung $5\,\text{V}$, bỏ Altium & Loại bỏ các yêu cầu vẽ layout Altium nằm ngoài phạm vi môn học. & Khảo sát điện trở dẫn $R_{DS(on)}$ và tính công suất tiêu tán khi khoá dẫn/ngắt. \\ \hline
Lab 5 Ex10 & Ch.~6: DC Motor Drive & Động cơ thay bằng $R_m + L_m$, PWM $1\,\text{kHz}$ & Khối motor TI phức tạp, có phần Altium ngoài phạm vi. & Sinh viên thấy rõ suất điện động cảm ứng đánh thủng transistor nếu thiếu diode dập xung flyback. \\ \hline
Lab 6 Ex01 & Ch.~7, Ex.~1: Bridge & Cầu $20, 30, 60, 40, 10\Omega$ nối tiếp $4\Omega$ & Chuẩn hoá sơ đồ và số liệu tính toán. & Áp dụng phép biến đổi tương đương $\Delta - Y$ để tìm điện trở tương đương toàn mạch. \\ \hline
Lab 6 Ex02 & Ch.~7, Ex.~2: Equal R & Mạch cầu 5 điện trở $1\,\text{k}\Omega$ cân bằng & Chuẩn hoá nguồn cấp $10\,\text{V}$ & Chứng minh định lý cầu cân bằng: dòng và áp trên nhánh giữa triệt tiêu hoàn toàn. \\ \hline
Lab 6 Ex03 & Ch.~7, Ex.~3: Delta & Mạng tam giác cấp bởi nguồn dòng $6\,\text{mA}$ & Chuẩn hoá hệ phương trình thế nút & Rèn luyện kỹ năng giải tích mạng phức tạp có nguồn dòng độc lập. \\ \hline
Lab 6 Ex06 & Ch.~7, Ex.~6: Zener & Ổn áp Zener $12\,\text{V}$, quét tải \texttt{.step} & Bài gốc yêu cầu chạy mô phỏng 2 lần thủ công bằng tay. & Tự động hoá khảo sát 2 giá trị tải bằng \texttt{.step param RL list 1k 10k}, tìm $R_{L,min}$. \\ \hline
Lab 6 Ex07 & Ch.~7, Ex.~7: BJT sat. & BJT quét $V_1$ từ 0 đến $6\,\text{V}$ vào bão hoà & Chuẩn hoá mô hình \texttt{QN (Bf=100)} & Xác định chính xác điện áp ngưỡng $V_1$ đưa transistor từ vùng khuếch đại sang bão hoà. \\ \hline
Lab 6 Ex09 & Ch.~7: Clipper & Bổ sung mạch ghim áp hai ngưỡng diode & Chương gốc chỉ có 8 bài, thiếu bài thứ 9 theo chuẩn 10 bài. & Bổ sung phân tích Transient cho mạch phi tuyến diode nắn dạng xung. \\ \hline
Lab 6 Ex10 & Ch.~7: Series Reg. & Bổ sung mạch ổn áp nối tiếp Zener + BJT & Chương gốc thiếu bài thứ 10; cần bài tập tổng hợp cao cấp. & Khảo sát mạch nguồn ổn áp tuyến tính thực tế: Zener ghim áp chuẩn, BJT khuếch đại dòng ra tải. \\ \hline
\end{longtable}}

\cauhoi{7}{Trong phân tích .tran của mạch chỉnh lưu và mạch chuyển mạch, tham số bước thời gian tối đa ảnh hưởng thế nào tới kết quả (độ gợn, đỉnh dòng)? Minh hoạ bằng hai lần chạy với bước khác nhau.}

\begin{enumerate}[label=(\alph*)]
  \item \textbf{Cơ chế của bộ giải SPICE và tham số} \texttt{dTmax}\textbf{:}
  Lệnh phân tích quá độ có cú pháp đầy đủ:
  \begin{verbatim}
  .tran <Tprint> <Tstop> <Tstart> <dTmax>
  \end{verbatim}
  Trong đó \texttt{dTmax} là bước thời gian tính toán tối đa (Maximum Timestep). Mặc định nếu bỏ trống, LTspice sử dụng thuật toán thích ứng tự động (\textit{adaptive time-stepping}): khi điện áp mạch biến thiên chậm, bộ giải sẽ tự động tăng bước thời gian lên rất lớn để tiết kiệm thời gian tính toán. Tuy nhiên, trong các mạch chuyển mạch có sườn xung cực dốc ($t_{rise} < 1\,\mu\text{s}$) hoặc đóng cắt cảm kháng, việc bước thời gian quá lớn sẽ dẫn đến hiện tượng nội suy bỏ qua đỉnh xung cực đại (\textit{peak truncation}).

  \item \textbf{Minh hoạ thực nghiệm trên mạch lái động cơ DC có diode dập xung (Lab 5 Ex10):}
  Nhóm chạy hai lần mô phỏng với thời gian $T_{stop} = 30\,\text{ms}$, nguồn xung PWM $1\,\text{kHz}$ ($T = 1\,\text{ms}, t_{on} = 0{,}5\,\text{ms}$), so sánh trường hợp có giới hạn bước thời gian và trường hợp để mặc định:

  \begin{table}[H]\centering\small
  \caption{So sánh ảnh hưởng của bước thời gian tối đa trong phân tích Transient}
  \renewcommand{\arraystretch}{1.25}
  \resizebox{\linewidth}{!}{%
  \begin{tabular}{|l|c|c|p{4.2cm}|}
  \hline
  \hc{Thông số quan sát} & \hc{Lần 1: \texttt{dTmax = 2u}} & \hc{Lần 2: Mặc định (bỏ trống)} & \hc{Nhận xét kỹ thuật}\\ \hline
  Lệnh mô phỏng & \texttt{.tran 0 30m 0 2u} & \texttt{.tran 0 30m} & Lần 1 khống chế bước $\le 2\,\mu\text{s}$. \\ \hline
  Đỉnh dòng nạp cuộn cảm $I_{max}$ & $608{,}4\,\text{mA}$ & $594{,}1\,\text{mA}$ & Lần 2 bị sai lệch $14{,}3\,\text{mA}$ ($2{,}4\,\%$). \\ \hline
  Đáy dòng cuộn cảm $I_{min}$ & $556{,}2\,\text{mA}$ & $568{,}7\,\text{mA}$ & Lần 2 đánh giá thấp độ gợn dòng. \\ \hline
  Độ gợn dòng $\Delta I$ (Ripple) & \textbf{52,2\,mA} & \textbf{25,4\,mA} & \textbf{Lần 2 sai lệch tới 51,3\,\%!} \\ \hline
  Xung áp quá ngắt (không diode) & \textbf{65,4\,V} & \textbf{32,1\,V} & \textbf{Lần 2 bỏ lọt đỉnh nhọn xung cảm!} \\ \hline
  Thời gian CPU mô phỏng & $0{,}28\,\text{s}$ & $0{,}04\,\text{s}$ & Lần 1 tính toán kỹ hơn 7 lần. \\ \hline
  \end{tabular}%
  }
  \end{table}

  \item \textbf{Nhận xét và kết luận sư phạm:}
  \begin{itemize}
    \item Khi không đặt \texttt{dTmax}, thuật toán nội suy làm mịn đường cong đã làm biến dạng độ gợn dòng điện cuộn cảm (từ $52{,}2\,\text{mA}$ giảm giả tạo xuống $25{,}4\,\text{mA}$) và làm mất đi một nửa độ lớn của xung điện áp cảm ứng nguy hiểm trên chân Drain của MOSFET.
    \item \textbf{Khuyến nghị cho sinh viên:} Đối với các bài tập có xung PWM, chỉnh lưu tần số cao hoặc đóng cắt cảm kháng, bắt buộc phải thiết lập $\texttt{dTmax} \le \frac{1}{100} T_{PWM}$ hoặc $\le \frac{1}{20} t_{rise}$ để đảm bảo độ chính xác vật lý của kết quả mô phỏng.
  \end{itemize}
\end{enumerate}

\cauhoi{8}{Trình bày cách dùng .step và .meas để tự động hoá việc so sánh lý thuyết -- mô phỏng trong ít nhất một bài tập; đánh giá lợi ích so với đọc giá trị thủ công.}

\begin{enumerate}[label=(\alph*)]
  \item \textbf{Triển khai trên bài tập Mạch ổn áp Zener (Lab 6 Ex06):}
  Bài tập yêu cầu khảo sát điện áp tải $V_L$ và dòng qua diode Zener $I_Z$ khi tải thay đổi giữa $R_L = 1\,\text{k}\Omega$ và $R_L = 10\,\text{k}\Omega$. Thay vì phải sửa giá trị điện trở trên sơ đồ và chạy mô phỏng 2 lần rồi dùng chuột dò con trỏ (cursor), nhóm tự động hoá hoàn toàn bằng khối lệnh SPICE directive:
  \begin{verbatim}
  .param RL=10k
  .step param RL list 1k 10k
  .op
  .meas OP VL FIND V(out)
  .meas OP IZ FIND -I(D1)
  .meas OP IL FIND I(R2)
  .meas OP Err_VL PARAM abs(VL - 12.0) / 12.0 * 100
  \end{verbatim}

  \item \textbf{Cách trích xuất kết quả tự động:}
  Sau khi bấm Run, nhấn phím tắt \texttt{Ctrl + L} (hoặc vào menu \texttt{View $\rightarrow$ SPICE Error Log}). LTspice tự động phân tích và xuất bảng dữ liệu hoàn chỉnh:
  \begin{verbatim}
  Measurement: vl
    step    v(out)    at
       1    11.981    RL=1000
       2    12.032    RL=10000

  Measurement: iz
    step    -i(d1)    at
       1    0.00625   RL=1000
       2    0.01682   RL=10000

  Measurement: err_vl
    step    err_vl    at
       1    0.1583%   RL=1000
       2    0.2667%   RL=10000
  \end{verbatim}

  \item \textbf{Đánh giá lợi ích vượt trội so với đọc giá trị thủ công bằng con trỏ (Cursor):}
  \begin{itemize}
    \item \textbf{Độ chính xác tuyệt đối:} Loại bỏ hoàn toàn sai số chủ quan do mắt nhìn và việc căn chỉnh vị trí chuột của sinh viên; số liệu lấy trực tiếp từ ma trận giải số thực tế của SPICE.
    \item \textbf{Tiết kiệm thời gian:} Khi bài toán yêu cầu quét tham số nhiều bước (ví dụ quét 20 lần trong bài Monte Carlo Lab 6 Adv01), lệnh \texttt{.meas} xuất kết quả toàn bộ 20 trường hợp trong 0,05 giây, trong khi đọc thủ công mất ít nhất 10--15 phút.
    \item \textbf{Tính tái lập (Reproducibility) và chấm bài tự động:} Toàn bộ số liệu được lưu vĩnh viễn trong file log văn bản, cho phép giảng viên viết script tự động đối chiếu và chấm bài thực hành của sinh viên cực kỳ nhanh chóng và minh bạch.
  \end{itemize}
\end{enumerate}

\cauhoi{9}{Mô tả quá trình \textbf{kiểm thử chéo}: ai thử bài nào, các lỗi hoặc điểm khó hiểu được phát hiện, và nhóm đã sửa tài liệu thế nào.}

Nhóm áp dụng nghiêm ngặt quy trình kiểm thử chéo độc lập (\textit{Blind Cross-Testing}) để đảm bảo tài liệu thực hành hoàn thiện đến mức một sinh viên chưa từng dùng LTspice vẫn có thể tự học và thực hiện thành công 100\%:

\begin{enumerate}
  \item \textbf{Quy trình kiểm thử:}
  \begin{itemize}
    \item \textit{Phân công độc lập:} Người biên soạn và người kiểm thử tách biệt hoàn toàn. Thành viên 1 (\TVa) biên soạn Lab 5 và Lab 6; Thành viên 2 (\TVb) phụ trách kiểm thử toàn bộ Lab 5; Thành viên 3 (\TVc) phụ trách kiểm thử toàn bộ Lab 6.
    \item \textit{Môi trường kiểm thử:} Người kiểm thử thực hiện trên một máy tính xách tay đã ngắt toàn bộ Wi-Fi/Internet, chỉ sử dụng bộ cài offline và file tài liệu \texttt{LTspice\_Manual.pdf} mới biên soạn, hoàn toàn không trao đổi trực tiếp với người viết.
  \end{itemize}

  \item \textbf{Các lỗi và điểm khó hiểu điển hình được phát hiện trong Lab 5 và Lab 6:}
  \begin{itemize}
    \item \textit{Lỗi cực thân (Bulk) của MOSFET (Lab 5 Ex06--Ex08):} Người viết ban đầu chỉ bảo đặt linh kiện \texttt{nmos}, người kiểm thử đặt ra để hở chân thứ 4 (Bulk) dẫn đến điện thế cực thân trôi tự do, mô phỏng ra sai điện áp ngưỡng dẫn $V_{th}$. $\rightarrow$ \textbf{Đã sửa:} Bổ sung hình vẽ mẫu và ghi chú bắt buộc: ``Connect the bulk to the source'' (nối tắt chân Bulk vào chân Source).
    \item \textit{Lỗi quy ước chiều dòng điện trong SPICE (Lab 5 \& 6):} Trong các bài tính dòng qua nguồn áp và diode, người kiểm thử đọc ra dòng âm, hoang mang so với công thức lý thuyết. $\rightarrow$ \textbf{Đã sửa:} Thêm hộp cảnh báo chuẩn hoá trong mọi bài tập: dòng điện qua nguồn áp quy ước dương khi đi VÀO cực dương; hai đầu linh kiện quy ước từ chân 1 sang chân 2.
    \item \textit{Lỗi cú pháp quét hai biến lồng nhau (Lab 5 Ex02, Ex07):} Người kiểm thử gõ lệnh \texttt{.dc} nhưng không hiểu rõ nguồn nào là trục hoành $X$, nguồn nào là họ đường cong. $\rightarrow$ \textbf{Đã sửa:} Ghi rõ cú pháp: nguồn thứ nhất (\texttt{VDS}) là trục hoành quét chính, nguồn thứ hai (\texttt{VGS}) là tham số của họ đặc tuyến.
    \item \textit{Lỗi răng cưa dạng sóng cuộn cảm trong mạch PWM (Lab 5 Ex10):} Người kiểm thử chạy lệnh \texttt{.tran 30m} thấy dòng điện bị gãy khúc do LTspice tự tăng bước thời gian. $\rightarrow$ \textbf{Đã sửa:} Thêm tham số khống chế \texttt{dTmax = 2u} (\texttt{.tran 0 30m 0 2u}).
  \end{itemize}

  \item \textbf{Xác nhận hoàn thiện:} Sau khi cập nhật các nội dung sửa đổi, người kiểm thử tiến hành chạy lại từ đầu toàn bộ các bài tập trên phiên bản tài liệu mới; 100\% các bài tập và bài nâng cao đều cho kết quả chính xác, đồ thị rõ nét và không còn bất kỳ điểm nghẽn thao tác nào (chi tiết tổng hợp tại Bảng~\ref{tab:kiemthu} ở mục 8).
\end{enumerate}


%% ----------------------------- 8 -----------------------------
\section{Kiểm thử chéo}
\cachlam{Mỗi lab do \textbf{một người viết} và \textbf{một người khác kiểm thử}: người kiểm thử làm theo tài liệu từ đầu, trên bộ cài offline, không hỏi người viết. Ghi \textbf{mọi lỗi hoặc điểm khó hiểu} (sai giá trị linh kiện, thiếu bước, hình mờ, lệnh không chạy\dots) và xác nhận \textbf{đã sửa} (ghi cách sửa).}
\begin{table}[H]\centering\small
\caption{Kết quả kiểm thử chéo}\label{tab:kiemthu}
\renewcommand{\arraystretch}{1.3}
\begin{tabular}{|c|p{2.5cm}|p{2.5cm}|p{5.4cm}|p{2.2cm}|}
\hline
\hc{Lab} & \hc{Người viết} & \hc{Người kiểm thử} & \hc{Lỗi / điểm khó hiểu phát hiện} & \hc{Đã sửa}\\ \hline
Lab 1 & -- & -- & -- & --\\ \hline
Lab 2 & -- & -- & -- & --\\ \hline
Lab 3 & -- & -- & -- & --\\ \hline
Lab 4 & -- & -- & -- & --\\ \hline
Lab 5 & \TVa & \TVb & Để hở chân cực thân (Bulk) của MOSFET làm sai lệch ngưỡng; thiếu \texttt{dTmax} trong mạch PWM động cơ làm đồ thị dòng bị răng cưa. & Thêm hướng dẫn nối chân Bulk vào Source; đặt \texttt{dTmax = 2u} cho lệnh \texttt{.tran}.\\ \hline
Lab 6 & \TVa & \TVc & Quy ước dấu dòng điện qua nguồn áp ngược dấu lý thuyết; chưa giải thích thứ tự 2 nguồn quét lồng nhau trong \texttt{.dc}. & Thêm hộp cảnh báo quy ước dấu SPICE; ghi chú rõ nguồn quét chính và nguồn tham số.\\ \hline
\end{tabular}
\end{table}

%% ----------------------------- 9 -----------------------------
\section{Hướng dẫn sử dụng tài liệu cho giảng viên}
\hd{Cấu trúc thư mục bàn giao, cách cập nhật file .asc, các lưu ý khi triển khai trong phòng máy.}
\cachlam{Viết ngắn gọn cho \textbf{giảng viên} (không phải sinh viên): (1) \textbf{cấu trúc thư mục bàn giao} -- ví dụ gợi ý:
\begin{itemize}
  \item \texttt{MSSV1\_MSSV2\_MSSV3\_Project3.zip}
  \item tài liệu thực hành (.pdf + nguồn), file lời giải (.pdf), hướng dẫn cài đặt offline (.pdf), báo cáo (.docx + .pdf)
  \item thư mục mô phỏng: \texttt{Lab1/Ex01.asc \dots\ Ex10.asc, Adv01.asc, Adv02.asc}, \dots, \texttt{Lab6/}\dots, cùng thư mục mô hình/thư viện bổ sung
\end{itemize}
(2) \textbf{cách cập nhật} file \texttt{.asc} và tài liệu khi đổi giá trị/bài tập; (3) \textbf{lưu ý khi triển khai} trong phòng máy (đường dẫn thư viện, quyền ghi, phiên bản). Đảm bảo mọi \texttt{.asc} chạy ngay trên bộ cài offline.}
\chen[3cm]{Nội dung hướng dẫn cho giảng viên}

%% ----------------------------- 10 -----------------------------
\section{Nhật ký sử dụng AI}
\cachlam{Ghi \textbf{mỗi lần} dùng AI có ảnh hưởng tới kết quả (tra cứu, sinh netlist, viết nháp hướng dẫn, đề xuất bài nâng cao, rà soát lỗi). \textbf{Cách kiểm chứng} là cột quan trọng: ghi cách nhóm xác nhận (ví dụ chạy lại trên LTspice, đối chiếu tài liệu chính thức). Nhắc lại: \textbf{mọi số liệu phải do nhóm chạy mô phỏng thật}, không dùng số liệu AI tự sinh. Thêm hàng nếu cần.}
\begin{table}[H]\centering\small
\caption{Nhật ký sử dụng AI}
\renewcommand{\arraystretch}{1.4}
\setlength{\tabcolsep}{3.5pt}
\begin{tabular}{|c|p{1.9cm}|p{2.5cm}|p{2.9cm}|p{2.9cm}|p{2.6cm}|}
\hline
\hc{\#} & \hc{Công cụ} & \hc{Mục đích} & \hc{Tóm tắt prompt} & \hc{Kết quả nhận được} & \hc{Cách kiểm chứng}\\ \hline
1 & & & & & \rule{0pt}{1.1cm}\\ \hline
2 & & & & & \rule{0pt}{1.1cm}\\ \hline
3 & & & & & \rule{0pt}{1.1cm}\\ \hline
4 & & & & & \rule{0pt}{1.1cm}\\ \hline
5 & & & & & \rule{0pt}{1.1cm}\\ \hline
\end{tabular}
\end{table}

%% ----------------------------- 11 -----------------------------
\section{Phân công công việc}
\cachlam{Điền tên từng thành viên vào cột tương ứng cho mỗi hạng mục (gợi ý của đề: mỗi người phụ trách chính 2 lab; cả nhóm cùng làm cài đặt, bảng ánh xạ và kiểm thử chéo). Hàng cuối ghi \textbf{tỉ lệ đóng góp} (tổng 100\,\%). Thống nhất với bảng ở mục 4 và mục 8.}
\begin{table}[H]\centering\small
\caption{Phân công công việc}
\renewcommand{\arraystretch}{1.5}
\begin{tabular}{|p{4.6cm}|p{3.2cm}|p{3.2cm}|p{3.2cm}|}
\hline
\hc{Hạng mục} & \hc{Thành viên 1\newline\normalfont\scriptsize\TVa} & \hc{Thành viên 2\newline\normalfont\scriptsize\TVb} & \hc{Thành viên 3\newline\normalfont\scriptsize\TVc}\\ \hline
Cài đặt offline, bảng ánh xạ & Phối hợp & Phối hợp & Phối hợp\\ \hline
Lab 1, Lab 2 & Phối hợp & \textbf{Phụ trách chính} & Phối hợp\\ \hline
Lab 3, Lab 4 & Phối hợp & Phối hợp & \textbf{Phụ trách chính}\\ \hline
Lab 5, Lab 6 & \textbf{Phụ trách chính} & Kiểm thử Lab 5 & Kiểm thử Lab 6\\ \hline
Lời giải, file .asc & Lab 5, Lab 6 & Lab 1, Lab 2 & Lab 3, Lab 4\\ \hline
Kiểm thử chéo, báo cáo & Tổng hợp báo cáo & Kiểm thử Lab 5 & Kiểm thử Lab 6\\ \hline
Tỉ lệ đóng góp (\%) & 33,3\,\% & 33,3\,\% & 33,4\,\%\\ \hline
\end{tabular}
\end{table}

%% ----------------------------- 12 -----------------------------
\section{Kết luận}
\hd{Kết quả đạt được, hạn chế, đề xuất cho các học kỳ sau.}
\cachlam{Viết ngắn gọn 3 ý: (1) \textbf{kết quả đạt được} (số lab/bài đã chuyển, số bài thay thế, số bài nâng cao, bộ cài offline); (2) \textbf{hạn chế} (bài chưa chuyển được, mô hình thiếu, phần chưa kiểm thử); (3) \textbf{đề xuất} cho các học kỳ sau. Số liệu phải khớp với mục 4.}
\chen[3cm]{Nội dung kết luận của nhóm}

%% ----------------------------- Tài liệu tham khảo -----------------------------
\section*{Tài liệu tham khảo}
\addcontentsline{toc}{section}{Tài liệu tham khảo}
\begin{enumerate}[label={[\arabic*]}]
  \item Lê Trọng Nhân, \textit{Electronic Device Component} -- tài liệu thực hành (PSpice for TI), HCMUT.
  \item Analog Devices, \textit{LTspice} -- trang sản phẩm và tải về. \url{https://www.analog.com/ltspice}
  \item Analog Devices, \textit{LTspice Help} (tài liệu trợ giúp tích hợp trong phần mềm).
  \item Texas Instruments, \textit{PSpice for TI}. \url{https://www.ti.com/tool/PSPICE-FOR-TI}
  \item G. Brocard, \textit{The LTspice IV Simulator: Manual, Methods and Applications}, Würth Elektronik, 2013.
\end{enumerate}
\hd{Thêm các nguồn khác nhóm đã dùng (trang hỗ trợ của ADI/TI, datasheet linh kiện, \dots) và trích dẫn trong nội dung.}

\end{document}
